\documentclass[9pt,a4paper]{extarticle}
\usepackage[margin=1.0cm,top=1.3cm,bottom=1.2cm,headsep=0.25cm]{geometry}
\usepackage[T1]{fontenc}
\usepackage{lmodern}
\usepackage{microtype}
\usepackage{amsmath,amssymb}
\usepackage[dvipsnames,table]{xcolor}
\usepackage{booktabs,longtable,array,tabularx}
\usepackage{enumitem}
\usepackage{fancyhdr}
\usepackage{listings}
\usepackage[colorlinks=true,linkcolor=NavyBlue,urlcolor=NavyBlue,citecolor=NavyBlue]{hyperref}

\setlist{itemsep=1pt,topsep=3pt,parsep=0pt}
\setlength{\parindent}{0pt}
\setlength{\parskip}{2pt}
\renewcommand{\arraystretch}{1.05}
\setcounter{tocdepth}{2}
\setcounter{secnumdepth}{2}
\emergencystretch=2em

% ---------- column helpers ----------
% widths are scaled by 1.12 to fill the 19 cm text width
\newcolumntype{L}[1]{>{\raggedright\arraybackslash}p{\dimexpr#1*112/100\relax}}
\newcolumntype{C}[1]{>{\centering\arraybackslash}p{\dimexpr#1*112/100\relax}}
\setlength{\tabcolsep}{4pt}
% light zebra striping on every table so rows are easy to follow
\rowcolors{2}{black!4}{white}

% ---------- colours ----------
\definecolor{pickc}{HTML}{1B7F3B}
\definecolor{avoidc}{HTML}{B3261E}
\definecolor{trapc}{HTML}{C26A00}
\definecolor{exc}{HTML}{1F5FA8}
\definecolor{keyc}{HTML}{555555}
\definecolor{decc}{HTML}{6A2C91}
\definecolor{hdrc}{HTML}{E8EEF7}

% ---------- boxes (no external packages: lrbox + fcolorbox) ----------
\newsavebox{\cboxbin}
\setlength{\fboxsep}{3pt}
\setlength{\fboxrule}{0.8pt}
\newenvironment{cbox}[3]{%
  \par\vspace{2pt}\noindent
  \def\cboxfc{#1}\def\cboxbg{#2}%
  \begin{lrbox}{\cboxbin}\begin{minipage}{\dimexpr\linewidth-2\fboxsep-2\fboxrule\relax}
  \setlength{\parskip}{3pt}\textbf{\textcolor{#1}{#3}}\ \ignorespaces
}{%
  \end{minipage}\end{lrbox}\fcolorbox{\cboxfc}{\cboxbg}{\usebox{\cboxbin}}\par\vspace{2pt}}

\newenvironment{pick}[1][Choose it when]{\begin{cbox}{pickc}{pickc!6}{#1:}}{\end{cbox}}
\newenvironment{avoid}[1][Do NOT choose it when]{\begin{cbox}{avoidc}{avoidc!5}{#1:}}{\end{cbox}}
\newenvironment{trap}[1][Exam trap]{\begin{cbox}{trapc}{trapc!7}{#1:}}{\end{cbox}}
\newenvironment{ex}[1][Worked example]{\begin{cbox}{exc}{exc!5}{#1.}}{\end{cbox}}
\newenvironment{key}[1][Key facts]{\begin{cbox}{keyc}{black!4}{#1.}}{\end{cbox}}
\newenvironment{dec}[1][Decision rule]{\begin{cbox}{decc}{decc!5}{#1:}}{\end{cbox}}
\definecolor{profc}{HTML}{00695C}
\newenvironment{prof}[1][From the professor]{\begin{cbox}{profc}{profc!6}{#1:}}{\end{cbox}}

% ---------- small helpers ----------
% source tag: which deck(s) to open in the exam
% slide-source tags are hidden in the exam print (slides are not allowed in the exam);
% restore the line below to show them again
% \newcommand{\src}[1]{\par\vspace{-2pt}\hfill{\footnotesize\itshape\textcolor{gray}{Slides: #1}}\par}
\newcommand{\src}[1]{\ignorespaces}
% material NOT stated in the slides (general knowledge / my derivation) -- verify before relying
\newcommand{\bs}{\textsuperscript{\textcolor{avoidc}{$\dagger$}}}
\newcommand{\cc}{\,cc}
\newcommand{\AMAT}{\mathrm{AMAT}}
\newcommand{\CPI}{\mathrm{CPI}}

\newcommand{\tblhead}[1]{\rowcolor{hdrc}#1}

% ---------- header/footer ----------
\pagestyle{fancy}
\fancyhf{}
\fancyhead[L]{\small\itshape Processor Design Exam Playbook}
\fancyfoot[C]{\small\thepage}
\renewcommand{\headrulewidth}{0.3pt}

\title{\LARGE\bfseries Processor Design Exam Playbook\\[2pt]
\normalsize Which architectural feature to choose, for which scenario, and why}
\author{}
\date{}
\usepackage{etoolbox}
\makeatletter\patchcmd{\@maketitle}{\vskip 2em}{\vskip 0pt}{}{}\makeatother

\begin{document}
\maketitle

\begin{center}\begin{minipage}{0.95\textwidth}\small
\noindent This is an open-book \emph{decision manual}. The exam gives a processor architecture, a workload, and a situation where the architecture is not working ideally; you diagnose why and redesign the whole architecture with justification. \textbf{How to use it in the exam:} (1) follow the 8-step method (\S\ref{sec:8step}) and the professor's pointer list (\S\ref{sec:pointers}); (2) underline clues and look them up in the master lookup (end of \S\ref{ch:method}); (3) find the symptom in \S\ref{sec:symptoms} to get the root cause and fix; (4) open the concept section it points to for detail and numbers; (5) check knock-on effects (\S\ref{sec:knockon}) and run the checklist (\S\ref{sec:checklist}). Scenarios 6--8 in \S\ref{ch:scen} are model answers in the exam format.
Every concept section follows the same pattern: a \textbf{decision map} (situation $\rightarrow$ feature), each feature with \textbf{what it fixes / what it costs / when it wins / when it loses}, \textbf{numbers you can quote}, and \textbf{exam traps}.

\medskip
\noindent\textbf{Conventions.}
A red dagger\bs{} marks something that is \emph{not} stated in the slides (general knowledge or arithmetic on slide numbers); prefer slide wording where both exist. Sections~\ref{ch:dlp} and~\ref{ch:pe} come from the textbook because the professor named those topics.
\end{minipage}\end{center}


\section{Read This First: The Exam and How to Attack It}\label{ch:method}

\subsection{What the exam looks like (from the professor)}\label{sec:examformat}
\begin{prof}
\begin{itemize}[leftmargin=*]
\item \textbf{Two design activities}, 30 marks for both, scaled down to 15. Subjective paper: no single right answer; marks go to \emph{reasoning}.
\item Each question gives you \textbf{(1) a processor architecture, (2) a workload, (3) a situation where the architecture is not working ideally}. Enough information is given; where it is not, \textbf{write your assumptions} and continue.
\item You must \textbf{understand the problem from the question itself} (the clues in the text give you $\sim$80\% of the answer), explain \textbf{why the problem is occurring}, and propose a solution.
\item Your solution must \textbf{adjust the whole architecture, not just one aspect}, and you must say \textbf{why your solution is better}, in detail, by connecting the concepts of the whole course.
\item Named topics: \textbf{SIMT / DLP} and how the processor exploits them; \textbf{multiple processing elements} and how the workload is distributed (in terms of timing); \textbf{benefits and demerits} of parallelism; \textbf{order of execution} identified from the workload; whether parallelism is present at all; \textbf{data, memory and communication} issues (parallelism may not benefit; will the memory hierarchy work?); \textbf{synchronisation}; \textbf{static vs.\ dynamic} decisions; \textbf{lab experience} (why performance was better or worse).
\end{itemize}
\end{prof}

\subsection{Professor's pointer $\rightarrow$ what you must write}\label{sec:pointers}
Every row is a mark you can collect. Tick them off on the question paper as you write.
\small
\begin{longtable}{L{3.6cm}L{7.4cm}L{4.6cm}}
\toprule
\tblhead{\textbf{He said\ldots} & \textbf{So in your answer, write\ldots} & \textbf{Where in this book}}\\
\midrule
\endhead
``Given: arch, workload, situation where it is not working ideally'' & It is a \emph{diagnosis} question. Start with ``Observed problem'' and ``Root cause''. Never jump straight to the fix. & \S\ref{ch:diag} (symptom $\rightarrow$ cause) \\
``Info will be given, mention assumptions'' & A numbered list: ``A1: L1 is private, 32\,KB, 64\,B lines (not stated).'' Use typical values; keep assumptions few and plausible. & Typical values: App.~\ref{sec:typical} \\
``Why is the problem occurring'' & Name the \emph{mechanism} (e.g.\ ``every increment invalidates the line in 7 other caches''), backed by a number or a clue quoted from the question. & \S\ref{sec:symptoms} \\
``Static / dynamic decision of the solution'' & For each choice: decided at compile/design time (static) or at run time by hardware/runtime (dynamic), and why that suits this workload. & \S\ref{sec:staticdyn} \\
``SIMT, DLP, how the processor exploits it'' & Is the same operation applied to many data? Then: lanes/warps, masking for branches, coalesced memory, latency hidden by switching warps. & \S\ref{ch:dlp} \\
``Multiple PEs, how the workload is distributed, timings'' & How you split the work (data/task/pipeline), static or dynamic assignment, a small timeline, $T_{par}=\max(\text{PE time})+T_{comm}+T_{sync}$. & \S\ref{ch:pe} \\
``Benefits and demerits of parallelism'' & Always both: speedup/throughput \emph{and} overheads (communication, sync, imbalance, coherence, power, contention). & \S\ref{sec:paraproscons} \\
``Order of execution, identify from workload'' & Draw the dependence graph; mark critical path, independent parts, loop-carried dependences, reductions, serial phases. & \S\ref{sec:order} \\
``Parallel present or not'' & Say explicitly which part is parallel (fraction $f$) and which is serial; apply Amdahl. & \S\ref{sec:order} \\
``Data, memory, communication; parallelism may not benefit'' & Check whether the workload is memory-, communication- or sync-bound \emph{before} adding PEs. & \S\ref{sec:paraproscons}, \S\ref{sec:roofline} \\
``If memory hierarchy will work or not'' & Compare working set to cache size, reuse vs.\ streaming, stride vs.\ block, sharing. Say yes/no and why. & \S\ref{sec:memworks} \\
``Why is your solution better'' & Compare with at least two alternatives: mechanism, a number, and the cost you accept. & Decision maps in every section \\
``Synchronisation'' & Which sync is needed (lock, barrier, atomic), its cost, how you reduce it (privatise, tree reduction, finer locks). & \S\ref{sec:sync} \\
``Answers generic but related to the question; connect every concept'' & Use course vocabulary, but every sentence must point to something in \emph{this} question. & \S\ref{sec:connect} \\
``Lab aspects, why performance is better'' & Quote one lab finding as evidence (``as observed in the coherence lab\ldots''). & \S\ref{sec:labs} \\
``Adjust the whole architecture, not just one aspect'' & Walk all layers: distribution $\rightarrow$ core $\rightarrow$ caches $\rightarrow$ memory $\rightarrow$ interconnect/coherence/sync $\rightarrow$ power, and show knock-on effects. & \S\ref{sec:layers}, \S\ref{sec:knockon} \\
\bottomrule
\end{longtable}
\normalsize

\subsection{The 8-step method (use it for both questions)}\label{sec:8step}
\begin{enumerate}[leftmargin=*]
\item \textbf{Extract the givens} into three short lists: \emph{Architecture} (PEs/cores, issue width, in-order/OoO, threads, caches, memory, interconnect); \emph{Workload} (operations, data size, access pattern, dependences, branches, threads); \emph{Symptom} (what is not ideal, with its number).
\item \textbf{State assumptions} for anything missing (numbered A1, A2, \ldots).
\item \textbf{Characterise the workload}: kind and amount of parallelism (ILP/DLP/TLP), order of execution (dependence graph, critical path, serial fraction), memory behaviour (working set, reuse, stride), sharing/communication, branches. Write ``parallelism present: yes/no, where''.
\item \textbf{Diagnose}: match the symptom to a root cause (\S\ref{ch:diag}). Confirm with one number (AMAT, CPI, Amdahl, bandwidth, timeline) and rule out one competing explanation.
\item \textbf{Solution across the whole architecture}: go layer by layer (\S\ref{sec:layers}); one fix per bottleneck; then check what each fix needs elsewhere (\S\ref{sec:knockon}).
\item \textbf{Static or dynamic?} For each choice, say which and why (\S\ref{sec:staticdyn}).
\item \textbf{Justify}: why better (mechanism + number), its cost/demerit, rejected alternative with reason.
\item \textbf{Expected outcome}: what improves, roughly how much, and where the bottleneck moves next (it always moves).
\end{enumerate}

\begin{key}[One-sentence justification pattern]
``Because the workload is \underline{X} (evidence: \underline{clue/number}), the bottleneck is \underline{Y}; therefore I choose \underline{feature}, which reduces \underline{CPI term / AMAT term / $T_{comm}$ / $T_{sync}$ / power} by \underline{mechanism}, at the cost of \underline{cost}; it is decided \underline{statically/dynamically} because \underline{reason}; I rejected \underline{alternative} because \underline{constraint}.''
\end{key}

\textbf{Answer layout (headings to copy):} \textbf{1.} Given \& assumptions $\cdot$ \textbf{2.} Workload analysis (parallelism, order of execution, memory, communication) $\cdot$ \textbf{3.} Why the architecture is not ideal (root cause + evidence) $\cdot$ \textbf{4.} Proposed architecture (table: layer $|$ change $|$ static/dynamic $|$ why) $\cdot$ \textbf{5.} Why this is better (vs.\ alternatives, numbers) $\cdot$ \textbf{6.} Demerits/costs and where the bottleneck moves. A small diagram (PEs + caches + interconnect, or a timeline) earns marks cheaply.
\textbf{Time per question:} $\approx$10\% reading and underlining clues, 15\% analysis and diagnosis, 45\% solution, 20\% justification, 10\% checklist (\S\ref{sec:checklist}).

\subsection{The layers you can adjust (``the whole architecture'')}\label{sec:layers}
If you have never designed a processor, this is the menu. A processor design \emph{is} these layers; a complete answer says something about every layer the question touches.
\small
\begin{longtable}{L{2.6cm}L{4.6cm}L{8.4cm}}
\toprule
\tblhead{\textbf{Layer} & \textbf{What it does} & \textbf{Knobs you can turn (section)}}\\
\midrule
\endhead
Software / work distribution & Splits the program and maps it onto hardware & Partitioning (data/task/pipeline), static vs.\ dynamic scheduling, granularity, data layout, privatisation, loop transformations (\S\ref{ch:pe}, \S\ref{ch:ilp}, \S\ref{ch:mem}) \\
Front end & Fetches/decodes instructions & Branch predictor, BTB, return stack, fetch/decode width (\S\ref{ch:pipe}) \\
Core back end & Executes instructions & Pipeline depth, issue width, in-order/OoO (scoreboard/Tomasulo/ROB), number and type of FUs, registers, VLIW vs.\ superscalar (\S\ref{ch:ilp}) \\
Threads per core & Hides stalls using other threads & Coarse/fine-grained MT, SMT (\S\ref{ch:tlp}) \\
Data-parallel units & Same op on many data & SIMD width, vector lanes, GPU/SIMT warps, masks (\S\ref{ch:dlp}) \\
Number/type of PEs & More workers & Cores, heterogeneous (big+small, CPU+GPU) (\S\ref{ch:tlp}, \S\ref{ch:pe}) \\
Caches & Keep data close & Size, associativity, block, levels, private/shared, victim cache, prefetch, non-blocking/MSHRs (\S\ref{ch:mem}) \\
Main memory & Bandwidth and latency & Banks/interleaving, width, channels, NUMA placement (\S\ref{ch:mem}) \\
Interconnect \& coherence & Lets PEs share data & Bus vs.\ network, snooping vs.\ directory, MSI/MESI (\S\ref{sec:coh}) \\
Synchronisation & Enforces ordering & Locks, barriers, atomics, reductions (\S\ref{sec:sync}) \\
Power & Energy budget & DVFS, power gating, core count vs.\ width (\S\ref{ch:fund}) \\
\bottomrule
\end{longtable}
\normalsize

\subsection{Master lookup: clue in the question $\rightarrow$ what it implies}
Underline these keywords in the question first. The columns give the default direction; later sections refine it.

\small
\begin{longtable}{L{3.6cm}L{6.5cm}L{5.8cm}}
\toprule
\tblhead{\textbf{If the scenario says\ldots} & \textbf{Lean towards\ldots} & \textbf{Lean away from\ldots}}\\
\midrule
\endfirsthead
\toprule
\tblhead{\textbf{If the scenario says\ldots} & \textbf{Lean towards\ldots} & \textbf{Lean away from\ldots}}\\
\midrule
\endhead
Battery / energy / thermal limit / PMD &
Simple or modest core, small caches, few levels, way prediction, DVFS, power gating, many small cores instead of one huge core &
Very wide issue, deep speculation, trace cache, aggressive prefetch (ILP techniques are \emph{energy-inefficient}) \\
Low cost / embedded &
Simple pipeline, static prediction / delay slot, small direct-mapped cache, VLIW if code is fixed &
Tomasulo+ROB, SMT, large L2/L3, large BTB \\
Server, many independent requests, database, web &
Multicore + multithreading (fine-grained/SMT), large shared LLC, banked memory, non-blocking caches, ECC &
Spending area on single-thread ILP \\
Single-thread latency, one big program &
Wide OoO speculative superscalar (Tomasulo+ROB), strong predictor + BTB + RAS, small fast L1 + big L2, prefetch &
Coarse-grained MT, VLIW (irregular code), more cores \\
Same operation on every element (image, matrix, ML, signal) &
DLP: SIMD/vector/GPU (SIMT); unit-stride/coalesced layout; banked/wide memory (\S\ref{ch:dlp}) &
More big OoO cores (area spent on control, not arithmetic) \\
\texttt{if} inside the data-parallel loop &
Masking/predication, group data so neighbours branch alike, keep divergent part small (\S\ref{sec:divergence}) &
Assuming full SIMT efficiency \\
Shared counters, histogram, global sum, ``all threads update'' &
Privatise per PE + tree reduction; pad to separate cache lines; atomics only at the end (\S\ref{sec:sync}) &
One global lock or one shared array updated by all \\
Locks, critical sections, barrier every iteration &
Finer-grained locks, test-and-test-and-set, fewer barriers (bigger chunks), lock-free reduction &
Coarse global lock; barrier per tiny step \\
Tasks of unequal or unknown size &
Dynamic scheduling (work queue / work stealing), smaller chunks (\S\ref{sec:distrib}) &
Static equal split (load imbalance) \\
Producer $\rightarrow$ consumer stages &
Pipeline parallelism across PEs, balance stage times, queues between stages &
Ignoring the slowest stage (it sets throughput) \\
Binary compatibility / legacy software &
Dynamically scheduled superscalar &
VLIW \\
Unpredictable latencies (cache misses) &
Dynamic scheduling, non-blocking cache, multithreading &
Static schedules with fixed latencies; VLIW lock-step \\
Deep pipeline / very high clock &
Strong branch prediction, BTB, small simple L1, pipelined cache &
Delayed branch, large associative L1 \\
High branch frequency &
2-bit $\rightarrow$ correlating/two-level/gshare, BTB, return stack &
Predict-not-taken alone, 1-bit predictor \\
Large working set, poor locality, pointer chasing &
Multithreading (fine/SMT) to hide latency, large LLC, victim cache &
Blocking/stream prefetch (need regular patterns) \\
Streaming data, used once, bandwidth-limited &
Wide/interleaved memory, more channels, stream prefetch, larger blocks; caches will \emph{not} help (no reuse) &
Bigger caches; more cores (they share the same bandwidth) \\
Many cores, data in far memory (NUMA) &
Place data near the PE that uses it, partition by owner, directory coherence &
Snooping bus beyond a few cores \\
Limited die area &
Multicore of smaller cores, SMT as cheap TLP (Power5: +24\% area) &
Wide issue \\
Real-time / predictability &
Static techniques, scratchpad or small caches\bs &
Heavy speculation, variable timing\bs \\
High availability (24/7) &
Redundancy, ECC; availability $=$ MTTF/(MTTF+MTTR) &
--- \\
\bottomrule
\end{longtable}
\normalsize

\begin{trap}[Common failure mode]
Choosing ``the most advanced thing'' for every box, or fixing only one layer. The exam rewards \emph{fit} and \emph{completeness}: a core fix that leaves the memory or synchronisation bottleneck in place gains little, and the professor said explicitly to adjust the whole architecture.
\end{trap}

\section{Diagnose: Why Is the Architecture Not Working Ideally?}\label{ch:diag}

\subsection{The bottleneck principle}
Performance is set by the \emph{tightest} of these limits; the ``non-ideal situation'' in the question is always one of them binding:
\[
T \;\ge\; \max\Big(\underbrace{\tfrac{\text{ops}}{\text{peak ops/s}}}_{\text{compute}},\;
\underbrace{\tfrac{\text{bytes}}{\text{bandwidth}}}_{\text{memory}},\;
\underbrace{\text{critical path}}_{\text{dependences}},\;
\underbrace{\text{serial part}}_{\text{Amdahl}}\Big) \;+\; T_{comm}+T_{sync}+T_{imbalance}
\]
Your diagnosis names the binding term and proves it with a number. Your solution shrinks that term, then you check which term binds next (\emph{the bottleneck moves}: Lab 8, \S\ref{sec:labs}).

\subsection{Symptom $\rightarrow$ root cause $\rightarrow$ fix $\rightarrow$ why it is better}\label{sec:symptoms}
\small
\begin{longtable}{L{3.0cm}L{4.1cm}L{4.4cm}L{4.0cm}}
\toprule
\tblhead{\textbf{Symptom / clue} & \textbf{Why it happens (root cause)} & \textbf{Fix (concepts to address)} & \textbf{Why better than alternatives}}\\
\midrule
\endfirsthead
\toprule
\tblhead{\textbf{Symptom / clue} & \textbf{Why it happens} & \textbf{Fix} & \textbf{Why better}}\\
\midrule
\endhead
Speedup flattens as PEs are added (e.g.\ 8 cores $\rightarrow$ 2.5$\times$) & (a) serial fraction (Amdahl); (b) shared memory bandwidth saturated; (c) sync/communication overhead grows with $p$. Tell apart: some PEs idle $\Rightarrow$ (a)/imbalance; all busy but stalled on memory $\Rightarrow$ (b); time in locks/barriers $\Rightarrow$ (c) & (a) parallelise or overlap the serial phase; (b) raise arithmetic intensity (blocking, fusion), more channels/banks, larger shared LLC; (c) privatise + reduce, fewer barriers, coarser tasks (\S\ref{ch:pe}) & Adding PEs attacks none of these; each fix targets the term that actually binds (Amdahl: $f=0.9 \Rightarrow$ max 10$\times$ however many PEs) \\
Gets \emph{slower} with more threads & Coherence ping-pong: true sharing (all write one counter) or false sharing (different variables, same line); lock contention; shared-cache thrashing & Privatise per thread, pad to line size, reduce at end; finer locks; partition or enlarge shared cache (\S\ref{sec:coh}, \S\ref{sec:sync}) & Removes invalidations per update instead of buying faster interconnect. Lab 5: IPC collapsed while instruction count stayed flat \\
Some PEs idle while others work & Load imbalance: static equal split of unequal tasks; or PEs wait on dependences & Dynamic scheduling (work queue / stealing), smaller chunks, cyclic distribution (\S\ref{sec:distrib}) & $T_{par}=\max$ PE time; balancing lowers the max. Dynamic adapts to unknown task sizes \\
Low IPC on a wide OoO core, few cache misses & Long dependence chain (loop-carried recurrence): only one op in flight; or one FU class saturated & Break the chain: multiple accumulators, unroll, reassociate; add units only of the saturated class; SMT to bring independent work (\S\ref{ch:ilp}) & Wider issue/larger window does nothing (Lab 8: 48/96-entry scheduler, 8-wide issue: same 123 cycles; Lab 10: recurrence gives 0.27 IPC vs.\ $\sim$3) \\
Front end cannot feed a wide back end & Decode width/complex instructions, taken branches, I-cache misses & Better predictor + BTB, uop cache/trace cache, decode-friendly code (\S\ref{ch:pipe}) & Lab 6: 6-instruction decode ceiling, not 32\,B/cycle, was the limit \\
Many mispredictions, deep pipeline & Data-dependent branches; penalty $\propto$ depth & Two-level/gshare predictor, predication (branch-free code), shallower pipe (\S\ref{ch:pipe}) & CPI$_{branch}=f_{br}\times$miss rate$\times$penalty; shrink the largest factor \\
High miss rate, working set $>$ cache & Capacity misses & Blocking/tiling, loop interchange, larger LLC, prefetch (\S\ref{ch:mem}) & Blocking fixes it in software at zero hardware cost \\
High miss rate, small working set, power-of-2 strides & Conflict misses (addresses alias to one set) & Associativity, victim cache, array padding (\S\ref{ch:mem}) & Lab 1: 100\% miss with a tiny working set; Lab 4: victim cache removed 61.6\% of cycles on a conflict pattern \\
Streaming, data used once, low FLOPs per byte & Memory-bandwidth bound; caches give only spatial locality & More channels/banks, wider memory, stream prefetch, fuse passes, smaller data types (\S\ref{sec:roofline}) & Bigger caches and more cores do nothing: the bytes still cross the memory bus \\
Pointer chasing / linked structures, latency bound & Each load depends on the previous: memory-level parallelism (MLP) $=1$ & Multithreading (fine/SMT/GPU warps), array-based layout, larger LLC, jump-pointer prefetch\bs & Lab 4: more MSHRs gave nothing for a dependent chase but 8$\times$ for 8 independent streams \\
Prefetcher reduces misses but run is slower & Prefetches late (untimely), evict useful data (pollution), steal bandwidth & Tune distance, throttle, use only for regular streams & Lab 4: stride prefetch cut misses 57.4\%$\rightarrow$23.2\% yet ran slower. Judge by time, not miss rate \\
Higher associativity: hit rate up but no gain / L2 miss rate rises & Hit time / clock stretched; misses move between levels & Keep L1 small and fast, associativity in L2; judge by AMAT & Lab 3: L2/L3 miss \emph{rate} rose while AMAT and DRAM traffic improved \\
SMT gives speedup $<1$ & Threads thrash shared caches/TLB & Disable SMT for such mixes, raise associativity, partition cache (\S\ref{ch:tlp}) & Power5: FP pairs had most conflicts and least gain \\
Fast on small input, collapses on large & Working set crosses a cache level or TLB reach & Blocking, larger LLC, larger pages\bs & Identify the size at which the cliff occurs and match tile size to it \\
SIMD/GPU far below peak & Divergence, uncoalesced/strided access, low occupancy, host-device transfer, tiny kernels & Data layout (SoA), coalescing, compaction/sorting, tile in shared memory, overlap transfers (\S\ref{ch:dlp}) & Bigger GPU does not help: waste scales with it \\
Barrier every iteration, PEs spend time waiting & Sync overhead plus imbalance paid at every barrier & Fewer barriers (larger chunks), point-to-point sync, asynchronous phases (\S\ref{sec:sync}) & One wait per phase instead of per step \\
Poor scaling beyond one socket (NUMA) & Remote memory accesses; data allocated on one node & Owner-computes partition, local allocation, directory coherence (\S\ref{sec:comm}) & Local access is several times faster than remote\bs \\
Throttling when all cores active & Power/thermal cap ($P\propto CV^2f$) & DVFS, fewer/smaller cores, power-gate idle units (\S\ref{ch:fund}) & Energy per task is the right metric \\
\bottomrule
\end{longtable}
\normalsize

\subsection{Will the memory hierarchy work for this workload?}\label{sec:memworks}
Rule: \textbf{caches work when data is reused before it is evicted} (reuse distance $<$ cache capacity) or when neighbouring words are used together (spatial locality). Check the access pattern:
\small
\begin{longtable}{L{4.1cm}L{2.8cm}L{8.7cm}}
\toprule
\tblhead{\textbf{Access pattern in the workload} & \textbf{Hierarchy works?} & \textbf{Reason and what to do}}\\
\midrule
\endhead
Working set fits in cache, reused (loops over small arrays) & Yes & Temporal locality. Keep the cache; nothing to add. \\
Sequential single pass (streaming) & Partly & Only spatial locality: miss rate $\approx 1/(\text{words per block})$. Bandwidth decides: prefetch, wide/interleaved memory; bigger cache useless. \\
Large matrices with reuse (matrix multiply) & No as written; yes after blocking & Capacity misses; blocking makes three $B\times B$ tiles fit, misses $2N^3\rightarrow 2N^3/B$. \\
Stride $\ge$ block size (column walk in row-major array) & No & One useful word per block. Interchange loops, change layout (array-of-structs $\leftrightarrow$ struct-of-arrays). \\
Power-of-2 strides, few arrays & No (conflicts) & Same set again and again: associativity, victim cache, pad arrays. \\
Random / pointer chasing & Poor & Little locality; hide latency with threads, or restructure data. \\
Read-shared by many PEs & Yes & Copies live in Shared state; a shared LLC holds one copy for all. \\
Written by many PEs (shared counters, flags) & Works \emph{against} you & Coherence misses (invalidations). Privatise, pad, reduce at end. \\
GPU, neighbouring threads use neighbouring data & Yes if coalesced & One wide transaction per warp; else up to 32 transactions. \\
\bottomrule
\end{longtable}
\normalsize

\subsubsection*{Roofline test: is it compute-bound or memory-bound?\bs}\label{sec:roofline}
\[
\text{Attainable FLOP/s}=\min\big(\text{peak FLOP/s},\; \text{AI}\times\text{peak memory bandwidth}\big),\qquad
\text{AI}=\frac{\text{FLOPs}}{\text{DRAM bytes}}
\]
\begin{ex}[Use it in one line]
$c[i]=a[i]+b[i]$ in doubles: 1 FLOP per 24 bytes, AI $=0.042$. With 100\,GB/s memory: at most $4.2$\,GFLOP/s, \emph{no matter how many cores, lanes or GPUs}. Conclusion to write: ``parallel compute will not benefit; the bandwidth term binds; raise AI (fuse loops, block for reuse) or buy bandwidth.'' Blocked matrix multiply has AI growing with tile size $B$, so it can become compute-bound.
\end{ex}

\subsection{Is parallelism present, and will it help? Benefits and demerits}\label{sec:paraproscons}
\small
\begin{longtable}{L{2.2cm}L{3.1cm}L{3.6cm}L{3.4cm}L{3.2cm}}
\toprule
\tblhead{\textbf{Kind} & \textbf{Exploits} & \textbf{Benefit} & \textbf{Demerits / overheads} & \textbf{Does NOT help when}}\\
\midrule
\endhead
Pipelining & Overlap of stages of successive instructions & Throughput up to 1 instr/cycle; higher clock & Hazards, branch/load penalties grow with depth; latch overhead & Dependences and branches dominate \\
ILP (superscalar, OoO) & Independent instructions in one thread & Speeds unmodified single-thread code; binary compatible & Area and power grow faster than width; saturates at 3--6 issue & Long dependence chains, unpredictable branches, memory-bound code \\
DLP (SIMD, vector, GPU/SIMT) & Same operation on many elements & One instruction controls many ops $\Rightarrow$ high ops per watt and per area; regular memory access & Divergence, alignment/coalescing, needs regular data; data transfer to accelerator & Branchy, irregular, pointer-based, or small data \\
TLP on one core (MT/SMT) & Independent threads sharing a core & Fills idle slots, hides memory latency cheaply & Threads compete for caches/FUs; single thread may slow & One thread only; cache-thrashing mixes \\
Multicore / many PEs & Independent threads/tasks & Near-linear throughput when tasks are independent; energy-efficient (smaller cores) & Communication, coherence, sync, load imbalance, shared bandwidth & Large serial fraction, fine-grained sharing, bandwidth-bound \\
Pipeline (task) parallelism & Producer$\rightarrow$consumer stages on different PEs & Throughput even when each item is sequential & Throughput $=$ slowest stage; buffers; fill/drain latency & Stages unbalanced \\
\bottomrule
\end{longtable}
\normalsize
\begin{dec}[Parallelism will NOT benefit when]
(1) the serial fraction is large (Amdahl); (2) the workload is memory-bandwidth bound (roofline); (3) communication/computation ratio is high (tasks too fine); (4) PEs must synchronise or share writable data frequently; (5) tasks are unbalanced and statically assigned; (6) the problem is too small to amortise start-up/transfer; (7) power cap forces lower clock per core. Say which of these apply \emph{before} proposing more PEs.
\end{dec}

\subsection{Static vs.\ dynamic: the cross-cutting decision}\label{sec:staticdyn}
\textbf{Static} = decided before run time (compiler, programmer, design time). \textbf{Dynamic} = decided at run time (hardware or runtime system).
\small
\begin{longtable}{L{2.6cm}L{3.0cm}L{3.2cm}L{3.3cm}L{3.3cm}}
\toprule
\tblhead{\textbf{Decision} & \textbf{Static form} & \textbf{Dynamic form} & \textbf{Static wins when} & \textbf{Dynamic wins when}}\\
\midrule
\endhead
Instruction scheduling & Compiler scheduling, unrolling, software pipelining, VLIW & Scoreboard, Tomasulo + ROB (OoO) & Latencies fixed and known, regular loops, power/area tight & Cache misses, aliasing, binary compatibility needed \\
Branch handling & Predict-not-taken, BTFNT, profile, delay slot & 2-bit BHT, correlating/gshare, BTB & Shallow pipe, loop-dominated, cost-sensitive & Deep pipe, data-dependent branches \\
Register naming & Compiler allocation & Hardware renaming & Enough architectural registers & WAR/WAW from few registers, OoO \\
Multiple issue & VLIW (compiler packs) & Superscalar (hardware picks) & Fixed product, strong compiler & Varying code, compatibility \\
Work to PEs & Block/cyclic partition fixed in code & Work queue, work stealing, guided chunks & Equal, predictable task costs; locality matters; zero overhead & Unequal or unknown costs, heterogeneous PEs \\
Prefetching & Compiler prefetch instructions & Hardware stride/stream prefetcher & Access pattern known at compile time & Pattern visible only at run time \\
On-chip data placement & Scratchpad / GPU shared memory (software managed) & Hardware cache & Access pattern known (tiles); predictable timing & Irregular access; ease of programming \\
Thread $\rightarrow$ core mapping & Pinning, owner-computes & OS/runtime migration & NUMA locality, cache affinity & Load varies over time \\
Power & Fixed V/f & DVFS, power gating & Constant load, real-time & Bursty load, battery \\
Cache sharing & Partitioned ways per core & Fully shared LLC & Isolation/QoS needed & Unequal demand between cores \\
GPU branch handling & Compiler predication (short branches) & Hardware masks/reconvergence & Short \texttt{if} bodies & Long divergent paths \\
\bottomrule
\end{longtable}
\normalsize
\begin{key}[Rule to write in every answer]
\textbf{Predictable and known before run time $\rightarrow$ static} (cheaper, lower power, deterministic, no run-time overhead). \textbf{Data-dependent, input-dependent, or must survive hardware changes $\rightarrow$ dynamic} (adapts, but costs hardware, energy and run-time overhead). The best answer is often \textbf{hybrid}: static coarse partition for locality + dynamic stealing for balance; compiler scheduling + OoO core.
\end{key}

\section{Fundamentals, Constraints and Trade-offs}\label{ch:fund}
This chapter holds the \emph{constraints} an exam scenario will give you (power, cost, availability, class of machine) and the quantitative tools to reason about them. Everything later is a way of improving one term of the equations here.

\subsection{Classes of computers: what each one optimises}
\small
\begin{tabularx}{\textwidth}{L{2.8cm}L{4.2cm}X}
\toprule
\tblhead{\textbf{Class} & \textbf{Emphasis (slides)} & \textbf{Design consequence}}\\
\midrule
Personal mobile device (phones, tablets) & Energy efficiency, real-time & Low-power cores, small/few cache levels, DVFS and power gating; throughput per joule beats peak speed. \\
Desktop & Price-performance & Single-thread performance still matters; balanced OoO core with moderate multicore. \\
Server & Availability, scalability, throughput & Multicore/multithreading, large caches, high memory bandwidth, ECC/redundancy. \\
Cluster / warehouse-scale (SaaS) & Availability, price-performance & Request-level parallelism; cost and energy per request dominate. \\
Supercomputer (sub-class) & FP performance, fast internal network & Wide FP units, high memory bandwidth, data-parallel structure. \\
Embedded & Price & Cheapest design that just meets the requirement; simple pipeline, small caches. \\
\bottomrule
\end{tabularx}
\normalsize
\src{ca01 ``Classes of Computers''}

\subsection{Kinds of parallelism and which hardware exploits them}
\begin{itemize}
\item \textbf{Application-level:} data-level parallelism (DLP: same operation on lots of data) and task-level parallelism (TLP: independent tasks).
\item \textbf{Architecture-level:} instruction-level parallelism (ILP: pipelining, superscalar, VLIW), vector/GPU (DLP), thread-level parallelism (multithreading, multicore), request-level parallelism (clusters/WSC).
\item \textbf{Flynn's taxonomy:} SISD; SIMD (vector architectures, multimedia extensions, GPUs); MISD (no commercial implementation); MIMD (tightly-coupled = multicore/shared memory; loosely-coupled = clusters).
\end{itemize}

\begin{dec}[Match workload to parallelism]
\emph{Same operation on huge regular arrays} $\rightarrow$ SIMD/vector/GPU (the slides only mention these as ``data-level parallelism''; modern graphics cards are SIMD). \emph{Independent threads/requests} $\rightarrow$ MIMD: multicore + multithreading. \emph{One sequential thread} $\rightarrow$ ILP hardware (pipeline, superscalar, speculation), but remember ILP saturates at roughly 3--6 issue.
\end{dec}
\src{ca01 ``Parallelism'', ``Flynn's Taxonomy''; ca08 ``General context: Multiprocessors''}


\subsection{The quantitative toolbox}
\subsubsection{Performance}
\[
\text{CPU time}=\text{IC}\times\CPI\times T_{clk},\qquad
\CPI=\frac{\sum_i \text{IC}_i\cdot\CPI_i}{\text{IC}}
\]
\begin{itemize}
\item \textbf{IC} (instruction count): ISA design and the compiler. \textbf{CPI}: microarchitecture (hazards, branches, memory stalls, issue width). \textbf{$T_{clk}$}: technology and pipeline depth.
\item Every technique in this manual attacks one of these factors, and usually \emph{trades} one against another (deeper pipeline: $T_{clk}\downarrow$ but $\CPI\uparrow$).
\item Speedup of X relative to Y $=$ execution time$_Y$ / execution time$_X$. Benchmarks: kernels, toy programs, synthetic, suites (SPEC, TPC-C).
\end{itemize}
\textbf{Amdahl's law:}\bs\ if a fraction $f$ of execution time is sped up by $s$, overall speedup $=1/\big((1-f)+f/s\big)$. Slide principle: \emph{focus on the common case}.
\src{ca01 ``Principles of Computer Design'', ``Measuring Performance''}

\subsubsection{Power and energy}
\begin{center}
\begin{tabular}{L{6.2cm}L{9.5cm}}
\toprule
\tblhead{\textbf{Quantity} & \textbf{Formula / statement (slides)}}\\
\midrule
Dynamic energy per switch & $\tfrac12\times C_{load}\times V^2$ \\
Dynamic power & $\tfrac12\times C_{load}\times V^2\times f_{switched}$ \\
Static power & $I_{static}\times V$; scales with number of transistors; reduce by \emph{power gating} \\
TDP (thermal design power) & Sustained power target for supply and cooling; lower than peak, higher than average \\
Better metric than power & Energy per task \\
\bottomrule
\end{tabular}
\end{center}
\begin{itemize}
\item \textbf{Reducing the clock rate reduces power but not energy per task} (slide). Lowering \emph{voltage} (DVFS) reduces energy because of the $V^2$ term.
\item Techniques the slides list: dynamic voltage-frequency scaling, low-power states for DRAM/disks, turning off cores, overclocking within thermal limits, power gating (static power).
\item Scale of the problem: Intel 80386 $\approx$2\,W; 3.3\,GHz Core i7 $\approx$130\,W from a 1.5\,cm$\times$1.5\,cm die, ``the limit of what can be cooled by air''.
\item If voltage scales with frequency, dynamic power falls roughly with $f^3$ (0.8$\times$ clock $\Rightarrow$ $\approx$0.51$\times$ power)\bs; this is why \emph{two slower cores can beat one fast core} on throughput per watt.
\end{itemize}

\begin{pick}[Energy-first design]
Choose DVFS/power gating, small caches, narrow issue and multicore of small cores whenever the scenario says battery, thermal envelope, or energy cost per request. Justify with $P\propto CV^2f$ and ``energy per task is the right metric''.
\end{pick}
\src{ca01 ``Power and Energy'', ``Dynamic Energy and Power'', ``Reducing Power'', ``Static Power''; ca06 ``Limits to ILP''}

\subsubsection{Cost}
\begin{itemize}
\item Cost falls along a \emph{learning curve} (yield improves). DRAM price tracks cost; microprocessor price depends on volume (10\% less per doubling of volume).
\item Die yield follows the Bose--Einstein form $\text{yield}=\text{wafer yield}\times(1+\text{defects/area}\times\text{die area})^{-N}$ with $N$ the process-complexity factor (slide: 11.5--15.5; defects 0.016--0.057/cm$^2$)\bs\ (formula is an image in the slide). \textbf{Consequence: bigger dies are disproportionately more expensive}, favouring several small cores over one huge core.
\end{itemize}
\src{ca01 ``Trends in Cost'', ``Integrated Circuit Cost''}

\subsubsection{Dependability}
\[
\text{MTBF}=\text{MTTF}+\text{MTTR},\qquad \text{Availability}=\frac{\text{MTTF}}{\text{MTBF}}
\]
For servers (availability is their emphasis) quote this and add redundancy and ECC (the slides mention error-correcting codes when discussing wide memory). \src{ca01 ``Dependability''; ca05 ``Wider Main Memory''}

\subsection{Technology trends that justify design choices}
\small
\begin{tabularx}{\textwidth}{L{4.2cm}X}
\toprule
\tblhead{\textbf{Trend (slide numbers)} & \textbf{What it implies for design}}\\
\midrule
Bandwidth improved 10{,}000--25{,}000$\times$ for processors but only 300--1200$\times$ for memory/disks; latency improved 30--80$\times$ vs.\ 6--8$\times$ &
Memory is the bottleneck (``memory wall''). Use caches, interleaving/wide memory, prefetching, and latency-hiding (OoO, non-blocking caches, multithreading). Bandwidth is easy to buy; latency is not. \\
Transistor density $+35\%$/yr; transistor speed scales, \textbf{wire delay does not} &
Huge centralised structures (many register ports, big CDB, wide issue logic) hurt clock rate. Favours smaller replicated cores and simpler structures. \\
Single-processor performance improvement ended in 2003; ILP cannot be leveraged further &
Move to explicit parallelism: DLP, TLP, RLP (``require explicit restructuring of the application''). Justifies multicore. \\
DRAM capacity growth slowing (25--40\%/yr); flash 15--20$\times$ cheaper per bit than DRAM; disk 15--25$\times$ cheaper than flash &
Deeper hierarchy: SRAM caches $\rightarrow$ DRAM $\rightarrow$ flash $\rightarrow$ disk, each level sized by cost per bit. \\
\bottomrule
\end{tabularx}
\normalsize
\src{ca01 ``Trends in Technology'', ``Bandwidth and Latency'', ``Transistors and Wires'', ``Current Trends in Architecture''}


\section{Pipeline, Long-Latency Units and Branch Handling}\label{ch:pipe}

\subsection{Decision map}
\small
\begin{tabularx}{\textwidth}{L{4.0cm}L{4.4cm}X}
\toprule
\tblhead{\textbf{Problem observed} & \textbf{Feature to apply} & \textbf{Why it fits}}\\
\midrule
Stalls from RAW hazards between dependent ALU ops & Forwarding/bypassing & Result is available at end of EX; forward instead of waiting for write-back. \\
Load followed by use & Compiler load-delay scheduling; forwarding from MEM & Load data is late; fill the slot(s) with independent instructions. \\
Control hazard, shallow pipe, predictable branches & Early PC evaluation in ID + static prediction / delay slot & Cheapest (no predictor memory), penalty already only 1 cycle. \\
Control hazard, deep pipe or high branch frequency & Dynamic prediction (BHT $\rightarrow$ correlating/two-level $\rightarrow$ gshare) + BTB & Penalty $=$ several cycles; accuracy matters more than anything. \\
Loop-heavy code & BTFNT or 2-bit counter & Back-edge is taken almost always; 2-bit avoids the double misprediction. \\
Target address needed early & BTB (and target buffer / branch folding) & Fetch target in IF with 0--1 bubble. \\
Returns / indirect jumps & Return-address stack & Targets vary at run time; stack predicts them. \\
Divide-heavy or multiply-heavy FP code & Pipelined (or replicated) FP units & Initiation interval 25 for an unpipelined divider blocks every other divide. \\
Want higher clock & Deeper pipeline (R4000 8-stage) & Shorter stages; but raises load/branch delay. \\
\bottomrule
\end{tabularx}
\normalsize

\subsection{Pipeline hazards: the three kinds and their standard cures}
\begin{center}
\small
\begin{tabular}{L{2.6cm}L{4.8cm}L{8cm}}
\toprule
\tblhead{\textbf{Hazard} & \textbf{Cause} & \textbf{Cure (slides)}}\\
\midrule
Structural & Two instructions need the same hardware in the same cycle & More hardware resources (replicate/pipeline units, more ports). \\
Data (RAW/WAR/WAW) & Dependence between instructions & Forwarding; compiler scheduling; (hardware) register renaming for WAR/WAW. \\
Control & Next PC unknown until branch resolves & Early evaluation of PC, delayed branch, prediction. \\
\bottomrule
\end{tabular}
\end{center}
\begin{itemize}
\item Pipelining improves instruction \emph{bandwidth} (throughput), not the latency of one instruction.
\item \textbf{Longer pipelines increase the impact of hazards}; interrupts, complex ISAs and FP all make pipelining harder.
\item Compilers reduce the cost of data and control hazards: load delay slots, branch delay slots, branch prediction hints.
\end{itemize}
\src{ca03\_noor ``Review: Summary of Pipelining Basics''}

\subsection{Long-latency functional units (FP pipeline)}
The slides' running example is a RISC-V/MIPS-style pipeline in which EX may repeat, and there are several functional units.

\begin{center}
\begin{tabular}{L{6cm}C{3cm}C{3.6cm}}
\toprule
\tblhead{\textbf{Functional unit} & \textbf{Latency} & \textbf{Initiation interval}}\\
\midrule
Integer ALU & 0 & 1 \\
Data memory (integer and FP loads) & 1 & 1 \\
FP add (4-stage pipelined) & 3 & 1 \\
FP/integer multiply (7-stage pipelined) & 6 & 1 \\
FP/integer divide (not pipelined) & 24 & 25 \\
\bottomrule
\end{tabular}
\end{center}
\textbf{Latency} $=$ intervening cycles between the producer and a consumer that avoids a RAW stall. \textbf{Initiation interval} $=$ cycles that must separate two operations of the \emph{same} type issued to the unit.

Latency table used in the compiler-scheduling examples (cycles between producer and consumer): FP ALU$\rightarrow$FP ALU 3; FP ALU$\rightarrow$store 2; load$\rightarrow$FP ALU 1; load$\rightarrow$store 0; integer$\rightarrow$integer 0.

\begin{itemize}
\item Longer pipeline $\Rightarrow$ more frequent RAW stalls (e.g.\ the 4th instruction after an FP add is the first that can use its result without a stall).
\item New hazards appear: \textbf{structural hazard at the register-file write port} (three instructions wanting write-back in the same cycle), \textbf{structural hazard on the unpipelined divider}, and \textbf{WAW hazards / out-of-order completion} because a short operation can finish before an earlier long one.
\end{itemize}

\begin{pick}[Pipelined vs.\ unpipelined divider]
Pipelined divider (or several dividers) when the code has frequent divides close together: with initiation interval 25, a second divide stalls for up to 24 cycles. An unpipelined divider is a cheap, acceptable choice when divides are rare, since it needs far less area and power. Adder and multiplier are worth pipelining almost always (initiation interval 1).
\end{pick}
\src{ca03\_noorNew ``RISC V FP Operation'', ``Latencies and initiation intervals'', ``A pipeline that supports multiple outstanding FP operations''; ca03\_noor ``FP Loop: Where are the Hazards?''}

\subsection{Deeper pipelines: the MIPS R4000 case study}
The R4000 splits the 5-stage integer pipeline into \textbf{8 stages}: IF, IS, RF, EX, DF, DS, TC, WB (cache access takes two stages each for instruction and data; TC is the tag check). The aim is a \textbf{higher clock rate}.

\begin{center}
\small
\begin{tabular}{L{5cm}L{10cm}}
\toprule
\tblhead{\textbf{Effect} & \textbf{Value / reason}}\\
\midrule
Load delay & \textbf{2 cycles} (the older deck states ``TWO cycle load latency''; the newer deck's slide title reads ``x1'' but its pipeline table shows two Stall cycles --- use 2). Data is available at end of DS and is bypassed; if TC reports a miss the pipeline backs up. \\
Branch delay & \textbf{3 cycles}: condition evaluated in EX; one delay slot plus two stalls; ``branch likely'' cancels the delay slot when not taken. \\
FP unit & Three units: divider, multiplier, adder; adder also used in the last step of multiply and divide. Stage kinds: U (unpack), S (operand shift), A (mantissa add), R (round), E (exception test), M, N (multiplier stages), D (divide stage, used 28 times in a row for divide). \\
\bottomrule
\end{tabular}
\end{center}

\begin{dec}[Deep vs.\ shallow pipeline]
\emph{Go deeper} when the target is a high clock rate and the code has good branch predictability and few back-to-back load-use pairs (scientific/loop code), \emph{and} you can afford a strong branch predictor. \emph{Stay shallow} for branchy, low-ILP, low-power or area-limited designs: every extra stage adds to the misprediction and load penalty (and, in general, to power\bs) (deeper pipelines also magnify branch cost: slide ``Branch Prediction Techniques'').
\end{dec}
\src{ca03\_noor ``Case Study: MIPS R4000''; ca03\_noorNew ``The MIPS R4000 Pipeline'' through ``The Floating-Point Pipeline''}

\subsection{The control hazard problem in numbers}
\subsubsection{Where the branch resolves determines the penalty (5-stage MIPS)}
\begin{center}
\begin{tabular}{L{8.6cm}C{3.2cm}}
\toprule
\tblhead{\textbf{Scheme} & \textbf{Stall cycles per branch}}\\
\midrule
Resolve at MEM, no forwarding (conservative stall) & 3 \\
Resolve at EX (with forwarding) & 2 \\
\textbf{Early evaluation: compare registers, compute target and update PC in ID} & \textbf{1} \\
\bottomrule
\end{tabular}
\end{center}
Early evaluation has a side-effect: an ALU instruction that produces the branch operand needs \textbf{1 extra stall}; a load that produces it needs \textbf{2 extra stalls}.

\[
\text{stall cycles per instruction (branches)} = f_{branch}\times \text{penalty};\qquad
\CPI = \CPI_{base}+f_{branch}\times(\text{mispredict rate})\times\text{penalty}\bs
\]
Even a 1-cycle penalty costs 10--30\% performance depending on branch frequency (slide). \textbf{Importance of prediction grows with (a) pipeline depth (penalty) and (b) branch frequency.}

\begin{ex}[Why prediction matters more in deep pipes]\bs
$f_{branch}=20\%$. Shallow pipe with early evaluation: always stall 1 cycle $\Rightarrow$ $+0.20$ CPI. Deep pipe, penalty 3: always stall $\Rightarrow +0.60$ CPI; with a 90\%-accurate predictor $\Rightarrow 0.20\times0.10\times3=+0.06$ CPI. The predictor recovers 90\% of the loss.
\end{ex}
\src{ca04 ``The Problem of Control Hazards'' through ``Branch Prediction Techniques''}

\subsection{Static branch prediction (decided at compile time)}
Static schemes need no run-time hardware, so they are used where branch behaviour is highly predictable at compile time, and also to assist dynamic predictors.

\small
\begin{longtable}{L{2.9cm}L{4.4cm}L{4.5cm}L{3.9cm}}
\toprule
\tblhead{\textbf{Scheme} & \textbf{How it works} & \textbf{Best scenario} & \textbf{Weakness}}\\
\midrule
\endhead
Predict not taken & Keep fetching fall-through; flush 1 instruction if taken (1-cycle penalty) & Branches mostly not taken (error checks, forward branches); simplest hardware. If 50\% not taken, halves control-hazard cost & Loops (back-edges are taken) \\
Predict taken & Fetch target once decoded & Only when the \emph{target is known before the outcome} (needs BTA in IF or a BTB). No advantage in the base MIPS pipeline & Needs extra target hardware \\
BTFNT (backward taken, forward not taken) & Direction of branch decides & Loop-dominated code: back-edge branches almost always taken & Irregular forward branches \\
Profile-driven & Compiler uses profile from earlier runs/hints & Stable behaviour across inputs (fixed embedded workloads) & Behaviour that changes with input \\
Delayed branch & Compiler fills slot(s) after branch with a useful instruction; executed whether or not taken & Shallow pipe with a single slot; compiler fills $\approx$50\% usefully & Deep pipelines (several slots, hard to fill); remaining slots become NOPs \\
\bottomrule
\end{longtable}
\normalsize

\subsubsection{Choosing where the delay-slot instruction comes from}
\begin{center}
\small
\begin{tabular}{L{2.6cm}L{6cm}L{6.2cm}}
\toprule
\tblhead{\textbf{Source} & \textbf{Use when} & \textbf{Legality}}\\
\midrule
From before & Always valid if an independent instruction exists before the branch & Safe: executes the same either way \\
From target & \textbf{Branch usually taken} (loop back-edges); target instruction must be copied & Must be harmless (e.g.\ writes an unused temporary) if branch falls through \\
From fall-through & \textbf{Branch usually not taken} (forward branches) & Must be harmless if branch is taken \\
\bottomrule
\end{tabular}
\end{center}
\textbf{Canceling (nullifying) branch:} the instruction encodes the predicted direction; if the prediction is wrong the delay-slot instruction becomes a NOP, so the compiler need not be conservative. MIPS \emph{branch-likely} is cancel-if-not-taken (slot filled from target) --- ideal for loop back-edges.
\src{ca04 ``Static Branch Prediction Techniques'' through ``Delayed Branch Technique'' (all pages)}

\subsection{Dynamic branch prediction (decided by hardware at run time)}
Two cooperating pieces: a \textbf{branch-outcome predictor} (taken?) and a \textbf{branch-target predictor/BTB} (where to?). The fetch unit uses them to pick the next PC in IF.

\subsubsection{Branch history table (BHT) and counters}
\begin{itemize}
\item \textbf{1-bit BHT:} table indexed by low-order bits of the branch PC; no tags (every access ``hits''; aliasing is tolerated because the entry is only a hint). Weakness: a loop branch (taken 9, not-taken 1) mispredicts \emph{twice} per loop execution (at exit, and on re-entry) $\Rightarrow$ only \textbf{80\% accuracy}.
\item \textbf{2-bit saturating counter:} must be wrong twice before the prediction flips; loop exit costs only one misprediction. Predict taken if counter $\ge$ half of max.
\item \textbf{n-bit counter:} generalisation; studies show 2-bit performs almost as well, so \emph{use 2-bit}.
\item A 4K-entry 2-bit BHT achieves 82--99\% accuracy (1--18\% misprediction) on SPEC89 and is almost as good as an infinite table; so beyond $\sim$4K entries, make the \emph{predictor smarter}, not bigger. Aliasing is cured by more rows or a hash.
\end{itemize}

\subsubsection{Correlating (two-level) predictors}
Idea: the outcome of recent \emph{other} branches helps predict this one (e.g.\ if the two preceding branches were both not taken, branch bb3 will be taken).
\begin{itemize}
\item \textbf{$(m,n)$ predictor:} the last $m$ branch outcomes (global history) select one of $2^m$ BHTs, each with $n$-bit counters. $(0,2)$ is the ordinary 2-bit predictor. $(2,2)$: four 2-bit BHTs selected by 2-bit history.
\item Index $=$ low-order PC bits concatenated with the $m$-bit history. Example: a $(2,2)$ predictor with 64 total entries uses a 6-bit index $=$ 2 history bits + 4 PC bits; each BHT has 16 entries.
\item A small $(2,2)$ predictor beats a plain 2-bit predictor of the same storage and often beats an unlimited-size 2-bit predictor.
\item \textbf{Two-level adaptive:} first level $=$ branch history register (BHR, a $k$-bit shift register of recent outcomes); second level $=$ pattern history table (PHT) of 2-bit counters indexed by the BHR.
\item \textbf{GAs:} PHT indexed by PC bits \emph{and} global history (local + global). \textbf{gshare:} PHT indexed by \textbf{XOR} of low PC bits with the BHR (local XOR global), which uses the table better than concatenation\bs.
\end{itemize}

\subsubsection{Branch target buffer (BTB) and fetch-unit support}
\begin{itemize}
\item BTB is a \emph{cache} accessed in IF with the PC of the instruction being fetched. Entry: \textbf{tag + predicted target (PC-relative) + optional prediction bits}. Only \textbf{taken} branches need an entry. The tag \emph{must} match (a wrong target cannot be used, unlike BHT aliasing).
\item If BTB hit and predicted taken: fetch from the predicted PC immediately. On misprediction: kill fetched instruction, restart at the other target, update/delete the entry. Correct prediction: no stalls.
\item \textbf{Target instruction buffers} store the target \emph{instruction}(s): faster, allow a larger BTB, and permit \textbf{branch folding} (0-cycle unconditional branches).
\item \textbf{Integrated instruction-fetch unit:} autonomous unit doing branch prediction, prefetch and instruction buffering to feed a wide back end.
\item \textbf{Return-address predictor:} small stack of recent return addresses, for indirect jumps/returns whose destination varies at run time.
\end{itemize}

\subsubsection{Which dynamic predictor for which scenario}
\small
\begin{longtable}{L{3.7cm}L{4.4cm}L{7.3cm}}
\toprule
\tblhead{\textbf{Scenario} & \textbf{Choose} & \textbf{Reason}}\\
\midrule
\endhead
Small/cheap core, loop code & 2-bit BHT & Fixes the loop double-miss at tiny cost (2 bits/entry, no tags). \\
Control-heavy integer code (compilers, interpreters), correlated if-chains & Correlating $(2,2)$ or two-level (GAs/gshare) & Uses other branches' outcomes; small table beats a much larger 2-bit table. \\
Same storage budget, want best accuracy & gshare & XOR indexing spreads history and PC across the table. \\
Wide-issue or deep pipeline & Two-level predictor + BTB + return-address stack, feeding \emph{speculation} & Cost of a miss is large; wide issue may need a branch every cycle. \\
Many function calls/returns & Return-address stack & Returns are indirect with run-time targets. \\
Fixed embedded workload, no predictor area & Static: profile-driven or BTFNT (+delay slot) & No hardware; behaviour known offline. \\
Branch target needed in IF for 0--1 bubble taken branches & BTB (+ target buffer for folding) & Predicts the \emph{address}, not just direction. \\
\bottomrule
\end{longtable}
\normalsize

\begin{trap}[Branch-prediction traps]
\begin{itemize}
\item 1-bit loop example: accuracy is 80\% for a 10-iteration loop (8 correct, 2 wrong), not 90\%.
\item ``Predict taken'' is useless unless the target address is available before the outcome (needs a BTB or early BTA).
\item A BHT has no tags (hint only); a BTB \emph{does} (an incorrect address cannot be used).
\item Prediction alone is not enough for wide-issue machines: it must be combined with \textbf{speculation} (Section~\ref{ch:ilp}).
\end{itemize}
\end{trap}
\src{ca04 ``Dynamic Branch Prediction Techniques'' through ``Branch Target Buffer'', ``High Performance Instruction Delivery''}


\section{Instruction-Level Parallelism: Scheduling, Speculation, Multiple Issue}\label{ch:ilp}

\subsection{Decision map}
\small
\begin{tabularx}{\textwidth}{L{4.2cm}L{4.3cm}X}
\toprule
\tblhead{\textbf{Scenario} & \textbf{Feature} & \textbf{Why}}\\
\midrule
Regular loops, independent iterations, fixed known latencies & Loop unrolling + static scheduling (+ software pipelining) & Compiler creates independent instructions to fill latency slots and removes loop overhead. \\
Unknown dependences at compile time (memory aliasing, cache-miss latencies) & Dynamic scheduling (Tomasulo / scoreboard) & Hardware reorders at run time; compiler need not know the microarchitecture. \\
Few architectural registers, lots of WAR/WAW & Tomasulo (register renaming in reservation stations) & Renames in hardware to a larger virtual register set. \\
Simplest possible out-of-order, few non-pipelined FUs & Scoreboard (CDC 6600 style) & Cheap centralised control; accepts stalls on WAR/WAW and no forwarding. \\
Need precise exceptions + branch speculation & Tomasulo $+$ reorder buffer (ROB) & Out-of-order execute, in-order commit; easy undo. \\
Need CPI $<1$ & Multiple issue: superscalar or VLIW & More than one instruction per cycle. \\
Binary compatibility across generations & Dynamically scheduled superscalar & Same binary, hardware finds ILP. \\
Fixed-product, compiler-controlled, power/area simplicity & VLIW (static scheduling) & No issue/hazard-detection hardware. \\
Wider than $\sim$4--6 issue is requested & Don't: go multicore/TLP instead & ILP saturates; energy inefficient; clock rate suffers. \\
\bottomrule
\end{tabularx}
\normalsize

\subsection{Dependences: what limits ILP}
\begin{center}
\small
\begin{tabular}{L{3.2cm}L{5.2cm}L{6.8cm}}
\toprule
\tblhead{\textbf{Type} & \textbf{Meaning} & \textbf{Can it be removed?}}\\
\midrule
True (data) dependence & Instruction $j$ uses a result produced by $i$ (RAW if a hazard) & \textbf{No.} It is the real data flow; instructions cannot execute in parallel. \\
Anti-dependence (name) & $j$ writes a register/memory location that earlier $i$ reads (WAR) & Yes: \emph{register renaming} (just a naming conflict). \\
Output dependence (name) & $i$ and $j$ write the same location (WAW); order must be kept & Yes: register renaming. \\
Control dependence & Instruction executes only if a branch goes a certain way & Relaxed by speculation/prediction; must preserve exception behaviour and data flow. \\
\bottomrule
\end{tabular}
\end{center}
\begin{itemize}
\item Register dependences are easy to detect (fixed names). \textbf{Memory dependences are hard}: does \texttt{100(R4)} equal \texttt{20(R6)}? Do different iterations touch the same address? (alias analysis).
\item \textbf{Control rules for code motion:} an instruction control-dependent on a branch cannot be moved before it; an instruction not control-dependent on it cannot be moved after it.
\item \textbf{Loop-carried dependence:} e.g.\ \texttt{A[i+1]=A[i]+C[i]} depends on the previous iteration, so iterations cannot run in parallel and the loop cannot be freely unrolled/parallelised. Safe to unroll when iterations are independent.
\end{itemize}
\src{ca03\_noor ``Compiler Perspectives on Code Movement'', ``When Safe to Unroll Loop?''; lec06 ``Review: Exceptions and Compiler Scheduling''}

\subsection{Static (compiler) scheduling}
\subsubsection{Loop unrolling: the slides' worked example}
Loop: \texttt{LD F0,0(R1); ADDD F4,F0,F2; SD 0(R1),F4; SUBI R1,R1,8; BNEZ R1,Loop} (+ delay slot).

\begin{center}
\small
\begin{tabular}{L{8.5cm}C{2.8cm}C{3.4cm}}
\toprule
\tblhead{\textbf{Version} & \textbf{Cycles} & \textbf{Cycles / iteration}}\\
\midrule
Original (stalls after LD, ADDD, and in branch slot) & 9 & 9 \\
Reordered: SUBI moved up, store address adjusted (\texttt{SD 8(R1),F4}) to fill delay slot & 6 & 6 \\
Unrolled $\times$4, straightforward (still stalls) & $15+4\times(1+2)=27$ & 6.75 \\
Unrolled $\times$4, scheduled with renamed registers (loads together, ADDDs together, stores together) & 14 & \textbf{3.5} \\
\bottomrule
\end{tabular}
\end{center}
\textbf{Why it works:} loop overhead (SUBI, BNEZ) is amortised; independent instructions from different iterations fill the 1-, 2- and 3-cycle latency gaps. \textbf{What it requires:} (a) \emph{renaming} (use F0/F4, F6/F8, \dots\ instead of reusing F0/F4, otherwise name dependences stall you), (b) knowledge that \texttt{0(R1)}, \texttt{-8(R1)}, \dots\ are different addresses so loads can move before stores, (c) no loop-carried dependence.

\subsubsection{Software pipelining and static renaming}
\begin{itemize}
\item \textbf{Loop unrolling:} multiple iterations in the loop body, all instructions from all iterations; needs compiler register renaming.
\item \textbf{Software pipelining:} one instruction from each of several iterations in the body; turns dependences into anti-dependences; often used for \emph{floating-point} with long latencies.
\item Static renaming limits: dependence is only resolved within a window, and \textbf{increases register pressure} for the compiler.
\end{itemize}

\begin{pick}[Static scheduling / unrolling]
Regular numeric loops with independent iterations and \emph{known fixed latencies}; in-order or VLIW hardware; embedded/DSP-style kernels where simple hardware and predictable timing matter. Best when many registers are available.
\end{pick}
\begin{avoid}
Code with unknown memory aliasing, pointer chasing, data-dependent branches, or variable latency from cache misses; or when the same binary must run efficiently on future hardware with different latencies.
\end{avoid}
\bigskip
\src{ca03\_noor ``FP Loop Hazards'' through ``Unrolled Loop That Minimizes Stalls''; lec06 ``Review''}

\subsection{Dynamic scheduling: why do it in hardware}
\begin{itemize}
\item Works when the dependence cannot be known at compile time; keeps the \emph{compiler simple}; code tuned for one machine runs well on another.
\item Key idea: \textbf{allow instructions behind a stalled one to proceed} (e.g.\ in \texttt{DIVD F0,F2,F4; ADDD F10,F0,F8; SUBD F12,F8,F14}, SUBD is independent and need not wait for the divide).
\item Consequence: out-of-order execution $\Rightarrow$ out-of-order completion $\Rightarrow$ \textbf{WAR and WAW hazards} become possible, and \textbf{precise exceptions} become hard.
\item Disadvantages: substantial hardware complexity; complicates exceptions.
\end{itemize}

\subsubsection{Scoreboard (CDC 6600, 1963)}
\textbf{In-order issue, out-of-order execution and commit.} Four stages replace ID/EX/WB:
\begin{enumerate}
\item \textbf{Issue:} functional unit free \emph{and} no active instruction with the same destination (no WAW); otherwise stall, and nothing behind issues.
\item \textbf{Read operands:} wait until both operands are available (RAW resolved here); no forwarding.
\item \textbf{Execute.}
\item \textbf{Write result:} stall if a WAR hazard exists with an earlier instruction that has not yet read its operand.
\end{enumerate}
Three tables: instruction status; functional-unit status (Busy, Op, Fi, Fj, Fk, Qj, Qk, Rj, Rk); register result status.

\textbf{Limitations (slide):} no forwarding hardware; limited to instructions in a basic block (small window, $\le5$ instructions in comparison table); few functional units (structural hazards, especially integer/load-store); no issue on structural hazard; \emph{waits} for WAR; \emph{prevents} WAW.

\subsubsection{Tomasulo's algorithm (IBM 360/91, 1966)}
Goal: high performance \emph{without special compilers}. The 360/91 had only \textbf{4 FP registers}, which blocked compiler scheduling; Tomasulo therefore \textbf{renamed registers in hardware}.
\begin{itemize}
\item \textbf{Distributed control:} \emph{reservation stations (RS)} sit with each functional unit and hold the op, operand values $V_j,V_k$, and tags $Q_j,Q_k$ (the RS that will produce the operand; $Q=0$ means ready). Register specifiers are replaced by values or RS tags $\Rightarrow$ \textbf{register renaming} that \textbf{eliminates WAR and WAW}. More RS than registers $\Rightarrow$ optimisations compilers cannot do.
\item \textbf{Common data bus (CDB):} results broadcast with their \emph{source} tag (``come-from'' bus: 64 data bits + 4 FU-id bits) to all RS, store buffers and the register file. Only the last writer updates the register.
\item Loads and stores are treated as functional units with buffers (load/store buffers hold address and data); memory order is kept through effective-address computation.
\item Three steps: \textbf{Issue} (if RS free, send operands/tags; else stall), \textbf{Execute} (when operands ready; watch the CDB otherwise), \textbf{Write result} (on CDB). No instruction starts until all earlier branches have resolved (non-speculative version).
\item \textbf{Overlaps loop iterations:} renaming gives different iterations different destinations (``dynamic loop unrolling''); RS let integer instructions run ahead past branches so FP ops beyond a basic block can be queued; also hides cache misses.
\end{itemize}

\begin{center}
\small
\begin{tabular}{L{3cm}L{6.2cm}L{6.2cm}}
\toprule
\tblhead{\textbf{Aspect} & \textbf{Tomasulo (IBM 360/91)} & \textbf{Scoreboard (CDC 6600)}}\\
\midrule
Functional units & Pipelined: 6 load, 3 store, 3 add, 2 mult/div & Multiple, not necessarily pipelined: 1 load/store, 1 add, 2 mult, 1 div \\
Window size & $\le 14$ instructions & $\le 5$ instructions \\
Structural hazard & No issue & Same \\
WAR & Renaming avoids stall & Stall completion \\
WAW & Renaming avoids stall & Stall issue \\
Result delivery & Broadcast from FU on CDB & Write/read registers (no forwarding) \\
Control & Distributed (reservation stations) & Centralised scoreboard \\
Worked example (6-instr.\ sequence) & Last instruction writes in cycle 57 & Same code finishes in cycle 62 \\
\bottomrule
\end{tabular}
\end{center}
Scoreboard is slower because of \textbf{structural hazards and lack of forwarding}.

\textbf{Tomasulo drawbacks (slide):} complexity; many associative stores (CDB) at high speed; \textbf{performance limited by the CDB} --- each CDB goes to many units (high capacitance, wiring density), and only \textbf{one FU can complete per cycle per CDB}; more CDBs need more associative logic; \textbf{imprecise interrupts} (fixed by the ROB).

\subsubsection{Hardware speculation: Tomasulo $+$ reorder buffer (ROB)}
Combines three ideas: \textbf{dynamic branch prediction} (choose what to run), \textbf{speculation} (execute before control dependences resolve, with ability to undo), \textbf{dynamic scheduling} (schedule across basic blocks). Core principle: \textbf{execute out of order, commit in order}; no irrevocable action (register/memory write, exception) before commit.

\begin{itemize}
\item \textbf{ROB:} circular buffer allocated in program order; entries hold instruction type (branch / store / register op), destination (register number or memory address), value, ready bit. It supplies operands between completion and commit and \emph{replaces the store buffer}.
\item \textbf{Four steps:} Issue (needs a free RS \emph{and} ROB slot) $\rightarrow$ Execute $\rightarrow$ Write result (CDB $\rightarrow$ ROB) $\rightarrow$ \textbf{Commit} (at ROB head: write register/memory; a mispredicted branch at the head \textbf{flushes the ROB} and restarts at the correct successor).
\item \textbf{Benefits:} \emph{precise interrupts} (exceptions recorded in the ROB, recognised only at commit; a speculative instruction's exception is flushed with it); \emph{early recovery} from mispredictions (clear entries after the branch, keep older ones); \emph{memory hazards}: stores update memory only at the head, eliminating WAR/WAW through memory; RAW through memory is kept by not letting a load proceed if an earlier store has a matching address.
\item \textbf{Costs:} large control complexity; associative network to find the latest version of a register (or a future file / register-status buffer); as many ROB ports as register-file ports; ROB sized to window.
\item \textbf{Design questions:} register renaming (extended physical registers) \emph{versus} ROB for holding values; \textbf{how much to speculate} --- speculation is hurt by exceptional events (cache miss, TLB miss) and speculating through multiple branches complicates recovery.
\end{itemize}

\begin{dec}[Scoreboard vs.\ Tomasulo vs.\ Tomasulo + ROB]
\begin{itemize}
\item \textbf{No dynamic scheduling} (in-order + compiler): cheapest/lowest power; embedded, predictable loops.
\item \textbf{Scoreboard:} want some out-of-order tolerance with minimal logic and few, non-pipelined units; accept WAR/WAW stalls.
\item \textbf{Tomasulo:} FP-heavy code with long latencies, \emph{few architectural registers}, or variable latencies (cache misses); you want renaming and loop overlap. Cap: one result per CDB per cycle.
\item \textbf{Tomasulo + ROB (speculative):} high-performance desktop/server cores needing precise exceptions and branch speculation, especially with multiple issue. Most area/power.
\end{itemize}
\end{dec}
\src{ca03\_noor (``HW Schemes'', ``Scoreboard'', ``Dynamic Scheduling - Tomasulo''); lec06 (whole deck: ``Tomasulo vs.\ Scoreboard'', ``Tomasulo Drawbacks'', ``Four Steps of Speculative Tomasulo''); ca04 ``Hardware-Based Speculation''}

\subsection{Multiple issue: superscalar vs.\ VLIW}
To get CPI $<1$ you must complete multiple instructions per clock. Three approaches: \emph{statically scheduled superscalar}, \emph{VLIW}, \emph{dynamically scheduled superscalar}.

\begin{center}
\small
\begin{tabular}{L{3cm}L{6.1cm}L{6.3cm}}
\toprule
\tblhead{\textbf{} & \textbf{Superscalar (dynamic)} & \textbf{VLIW (static)}}\\
\midrule
Who finds parallelism & Hardware at run time (RAW: forwarding; WAR/WAW: renaming) & Compiler at compile time: bundles of $K$ independent operations \\
Code size & Smaller & Larger: NOPs when bundles cannot be filled; unrolled loops \\
Binary compatibility & \textbf{Yes} across hardware generations & \textbf{Weak}: changing number of units ($K\to K+5$) or latencies breaks the schedule; fix = emulator, which hurts performance \\
Hardware & Complex issue/decode, dependence checking, renaming & \textbf{Simple decode/issue}, no interlock hardware (compiler checks) \\
Registers/ports & --- & More registers; simplified; may need multiple register files; VLIW example needs 7 read/3 write int and 5 read/3 write FP ports \\
Stalls & Independent instructions continue & \textbf{Lock-step}: one hazard stalls the whole bundle \\
Examples & Pentium 4, Athlon 64, Power5 & Itanium 2 (IA-64/EPIC) \\
\bottomrule
\end{tabular}
\end{center}
Example VLIW word: one integer op (or branch), two independent FP ops, two independent memory references. There must be enough parallelism to fill the slots.

\textbf{Dynamic multiple issue with speculation} (modern microarchitecture): limit instruction classes per bundle (e.g.\ one FP, one int, one load, one store); examine dependences among the bundle's instructions and encode them in the reservation stations; need multiple completion/commit (multiple CDBs). \textbf{Issue logic becomes the bottleneck.}

\begin{pick}[Superscalar (dynamic)]
General-purpose products with many software generations, legacy binaries, irregular integer code and unpredictable latencies.
\end{pick}
\begin{pick}[VLIW]
Embedded/DSP or fixed-function products where the compiler is controlled, the code is regular (loops with abundant ILP), the hardware must be simple, and binary compatibility across generations is \emph{not} needed.
\end{pick}
\src{ca03\_noor ``Static Instruction Scheduling'' through ``Dynamic Instruction Scheduling''; ca04 ``Multiple Issue and Static Scheduling'' through ``Superscalar v.\ VLIW''}

\subsection{Limits of ILP (use these to argue \emph{against} very wide issue)}
\subsubsection{Practical limits}
\begin{itemize}
\item \textbf{1 branch in 5 instructions:} how do you keep a 5-wide machine busy?
\item You need roughly \textbf{pipeline depth $\times$ number of FUs} independent instructions (e.g.\ $5\times4=15$--$20$) to keep the machine busy.
\item Hardware: \emph{easy} --- more instruction bandwidth, duplicate FUs; \emph{hard} --- more register-file ports; \emph{harder} --- more memory ports; wide decode hurts clock rate/pipeline depth.
\item Practical window sizes 32--126; larger windows are impractical (storage, comparisons, issue rate). Real processors are limited by number of FUs, buses, register ports.
\item \textbf{Doubling issue from 3--6 to 6--12} would need 3--4 data memory accesses/cycle, 2--3 branches resolved/cycle, renaming/accessing $>20$ registers/cycle, fetching 12--24 instructions/cycle $\Rightarrow$ the \textbf{clock rate suffers} (the widest-issue processor, Itanium 2, has the slowest clock and the highest power).
\item \textbf{Energy:} transistors switching $\propto$ peak issue rate, performance $\propto$ sustained rate; the gap grows, so \emph{multiple-issue techniques are energy-inefficient}.
\end{itemize}

\subsubsection{The ideal-processor study}
Assumptions: infinite renaming registers (no WAR/WAW), perfect branch and jump prediction, perfect caches (all loads/stores 1 cycle), perfect memory-address alias analysis, unbounded window and issue. ILP is then limited \emph{only by true data dependences}. Real model comparison: IBM Power5 issues 4/clk, 48 int + 40 FP rename registers, 2--6\% misprediction, 64\,KB I / 32\,KB D / 1.92\,MB L2 / 36\,MB L3.

\subsubsection{Head-to-head of four real ILP designs (Pentium 4 SMT, Athlon 64, Power5 SMT, Itanium 2 VLIW)}

\begin{center}
\small
\begin{tabular}{L{3.6cm}C{2.3cm}C{2.3cm}C{2.3cm}C{2.3cm}}
\toprule
\tblhead{\textbf{Efficiency rank (1=best)} & \textbf{Itanium 2} & \textbf{Pentium 4} & \textbf{Athlon 64} & \textbf{Power5}}\\
\midrule
SPECInt / transistor & 4 & 2 & 1 & 3 \\
SPECFP / transistor & 4 & 2 & 1 & 3 \\
SPECInt / area & 4 & 2 & 1 & 3 \\
SPECFP / area & 4 & 2 & 1 & 3 \\
SPECInt / Watt & 4 & 3 & 1 & 2 \\
SPECFP / Watt & 2 & 4 & 3 & 1 \\
\bottomrule
\end{tabular}
\end{center}
\textbf{Take-aways (slide):} no single leader; Athlon leads SPECInt; Itanium 2 and Power5 dominate on SPECFP; Itanium 2 is the least efficient on everything except SPECFP/Watt; Athlon and Pentium 4 use transistors and area well; Power5 is the best energy user on SPECFP and tied on SPECInt. Itanium (VLIW/EPIC) was ``not a significant breakthrough in scaling ILP''.

\begin{dec}[How wide should the core be?]
Quote the slide conclusion: \textbf{3--6 issue is the practical limit}. Choose a core width of about 4 for a general-purpose design; go narrower for energy/area-limited designs; beyond that, spend transistors on \emph{TLP} (multithreading, more cores). Right ILP/TLP balance is ``decided in the marketplace'': servers favour TLP, desktops may still prize single-thread speed.
\end{dec}
\src{ca04 ``Limits to Multi-Issue Machines'', ``Studies of the Limitations of ILP''; ca06 (whole deck)}



\section{Memory Hierarchy: Caches, Main Memory, Virtual Memory}\label{ch:mem}

\subsection{The framework: AMAT and the four levers}
\[
\AMAT = t_{hit}+\text{miss rate}\times\text{miss penalty},\qquad
\CPI_{total}=\CPI_{base}+\frac{\text{mem.\ refs}}{\text{instr.}}\times\text{miss rate}\times\text{miss penalty}\bs
\]
(The slides use exactly this CPI form in the memory-width example: $\CPI=2+1.2\times3\%\times64=4.3$.)

\begin{center}
\small
\begin{tabular}{L{3.8cm}L{11.6cm}}
\toprule
\tblhead{\textbf{Lever} & \textbf{Techniques (slides)}}\\
\midrule
Reduce \textbf{miss penalty} & Multilevel caches (the most important), critical word first \& early restart, read-miss priority over writes, merging write buffer, victim caches \\
Reduce \textbf{miss rate} & Larger blocks, larger caches, higher associativity, way prediction, pseudo-associative caches, compiler optimisations \\
Reduce penalty/rate via \textbf{parallelism} & Non-blocking caches, hardware prefetch, compiler-controlled prefetch \\
Reduce \textbf{hit time} & Small \& simple caches, avoid address translation in indexing, pipelined cache access, trace caches \\
\bottomrule
\end{tabular}
\end{center}
\textbf{Why hit time matters most:} it affects the clock cycle, so a fast hit helps \emph{everything} and its importance is ``multiplied'' beyond the AMAT formula; also miss penalty (in cycles) depends on clock cycle time.

\subsection{Step 1: diagnose the miss type (the ``3 C's'' + coherence)}
\small
\begin{longtable}{L{2.6cm}L{4.4cm}L{4.4cm}L{4.0cm}}
\toprule
\tblhead{\textbf{Miss type} & \textbf{Definition} & \textbf{Reduced by} & \textbf{Does NOT fix}}\\
\midrule
\endhead
Compulsory (cold) & First access to a block; measured by an infinite cache; independent of cache size & Larger blocks (spatial locality), prefetching & Bigger cache or more associativity \\
Capacity & Cache cannot hold all blocks needed; measured by fully associative cache of same size; falls as cache grows & Larger cache; blocking/tiling (shrinks working set); better temporal locality & Associativity \\
Conflict (collision) & Too many blocks map to the same set; vanishes with full associativity; falls as associativity rises & Higher associativity, victim cache, pseudo-associative, way-predicted set-assoc., code/data layout & Compulsory/capacity \\
Coherence (multiprocessor) & Block invalidated by another processor's write. \textbf{True sharing:} the same word is written by another processor. \textbf{False sharing:} different words that happen to be in the same cache line & Data layout (separate independently-updated data into different lines)\bs; smaller blocks\bs & Bigger caches \\
\bottomrule
\end{longtable}
\normalsize
\textbf{Slide examples:} 32\,B lines, array of 32 doubles: first touch of \texttt{array[0]} is a cold miss; prefetching \texttt{array[1..3]} turns later accesses into hits. In a direct-mapped cache, \texttt{array[0]}$\to$\texttt{array[16]}$\to$\texttt{array[0]} conflicts even though the cache is not full. In a multiprocessor, P2 writing \texttt{array[1]} invalidates P1's line holding \texttt{array[0]}: false-sharing miss.
\src{ca05 ``Classify Cache Misses --- 3 C's''; Cache-Misses (``Taxonomy of Cache Misses'')}

\subsection{Reducing miss penalty}
\subsubsection{Multilevel caches (``probably the best miss-penalty reduction'')}
\[
\AMAT=t_{hit,L1}+m_{L1}\big(t_{hit,L2}+m_{L2}^{local}\times P_{L2}\big),\quad
\text{global miss rate}=m_{L1}\times m_{L2}^{local}
\]
\begin{itemize}
\item \textbf{Local} miss rate $=$ misses in this cache / accesses \emph{to this cache}. \textbf{Global} $=$ misses / accesses generated \emph{by the CPU}. \textbf{Global miss rate is what matters} when judging L2/L3.
\item L1 capacity \emph{and} conflict misses are served by L2, avoiding main memory. A big L2 ($\gg$L1) gives a global miss rate close to a single-level cache of that size; the slide's conclusion is that L2 \emph{cache size} is what matters.
\item \textbf{L2 associativity:} direct-mapped hit 10\,cc, local miss 25\%; 2-way 10.1$\to$10 or 11\,cc, local miss 20\%; $P_{L2}=100$\,cc: DM $=10+0.25\times100=35$; 2-way $=10+0.20\times100=30$ (31 if the hit costs 11). \textbf{Higher associativity pays off at L2} because the extra hit time is tiny compared to the miss penalty.
\end{itemize}
\begin{ex}
1000 refs, 40 L1 misses, 20 L2 misses. $m_{L1}=4\%$; local $m_{L2}=20/40=50\%$; global $=2\%$. With $t_{L1}=1$, $t_{L2}=10$, $P_{L2}=100$, 1.5 refs/instr: $\AMAT=1+0.04(10+0.5\times100)=3.4$\,cc; stall cycles/instr $=(3.4-1)\times1.5=3.6$.
\end{ex}
\begin{pick}
Almost always. Essential when the processor is fast and memory is far (servers, desktops: L1+L2+L3). PMDs use two levels (64\,KB L1 + 256\,KB L2) to save energy/area. Hardware complexity 2 (costly if block sizes differ between levels).
\end{pick}

\subsubsection{Critical word first and early restart}
Do not wait for the whole block. \textbf{Critical word first:} fetch the missed word first and send it to the CPU, then fill the rest of the block. \textbf{Early restart:} fetch in normal order but restart the CPU the moment the requested word arrives. \textbf{Value depends on block size and spatial locality:} useful mainly for \emph{large blocks}; if the next access is the next sequential word in the block, the benefit shrinks.

\subsubsection{Read-miss priority over writes; merging write buffer}
\begin{itemize}
\item \textbf{Write-through with a write buffer:} a read miss may need data still sitting in the buffer (RAW through memory). Rather than waiting for the buffer to empty (increased read-miss penalty; one old MIPS 1000 design with a 4-word buffer saw about 50\% worse), \textbf{check the buffer for an address match and let the read go ahead if there is none}.
\item \textbf{Write-back:} on a read miss that replaces a dirty block, copy the dirty block to a buffer, do the read first, then write the dirty block: the CPU restarts sooner.
\item \textbf{Merging write buffer:} an entry holds several words but a write is usually one word. Merge a new write into an existing valid entry with a matching address: fewer full-buffer stalls and faster multiword writes. Mainly used with write-through caches.
\end{itemize}

\subsubsection{Victim cache}
A small \textbf{fully-associative} cache between the cache and its refill path holding only blocks \emph{just evicted}. On a miss, check it before going to the next level; if present, \textbf{swap} victim and cache block. No penalty increase because both are probed together.
\begin{itemize}
\item Jouppi: removes \textbf{20--95\% of the misses} of a \textbf{4\,KB direct-mapped} cache with 1--5 victim entries. AMD Athlon has 8 entries.
\end{itemize}
\begin{pick}
A direct-mapped (small, fast, simple) L1 suffering \emph{conflict} misses --- you keep the hit time of direct-mapped and recover much of 2-way's miss rate. Hardware complexity 2.
\end{pick}

\subsection{Reducing miss rate}
\subsubsection{Larger block size}
Reduces compulsory misses (spatial locality) but \textbf{raises miss penalty} and, in small caches, conflict and capacity misses. Pick the block that minimises \textbf{AMAT}, not miss rate.

\begin{center}
\small
\begin{tabular}{C{1.9cm}C{2.0cm}C{1.9cm}C{1.9cm}C{1.9cm}C{1.9cm}}
\toprule
\tblhead{\textbf{Block (B)} & \textbf{Penalty (cc)} & \textbf{4K} & \textbf{16K} & \textbf{64K} & \textbf{256K}}\\
\midrule
16 & 82 & 8.027 & 4.231 & 2.673 & 1.894 \\
32 & 84 & \textbf{7.082} & 3.411 & 2.134 & 1.588 \\
64 & 88 & 7.160 & \textbf{3.323} & \textbf{1.933} & \textbf{1.449} \\
128 & 96 & 8.469 & 3.659 & 1.979 & 1.470 \\
256 & 112 & 11.651 & 4.685 & 2.288 & 1.549 \\
\bottomrule
\end{tabular}
\end{center}
AMAT (cc) vs.\ block size; memory model: 80\,cc overhead + 2\,cc per 16\,B. \textbf{32--64\,B blocks dominate}; a 256\,B block only wins once the cache is 256\,KB (miss rate vs.\ block table: 4\,KB cache 256\,B block 9.51\% vs.\ 32\,B 7.24\%). Pentium 4's L2 uses 128\,B blocks.

\begin{dec}[Block size]
Small L1 $\rightarrow$ 32--64\,B. Large L2/L3 $\rightarrow$ 64--128\,B. Higher memory latency or bandwidth (overhead dominates the penalty) $\rightarrow$ somewhat larger blocks\bs. Small cache with 128/256\,B blocks $\rightarrow$ wrong (conflict and capacity misses rise).
\end{dec}

\subsubsection{Larger caches and higher associativity}
\begin{itemize}
\item \textbf{Larger cache:} helps capacity and conflict misses; costs hit time, area, power; popular for off-chip/L2.
\item \textbf{Associativity:} 8-way is nearly as good as fully associative. \textbf{2:1 rule:} 2-way of size $N/2$ $\approx$ direct-mapped of size $N$ (for caches $<$128\,KB). Cost: longer hit time $\Rightarrow$ clock cycle $\times1.10$ (2-way), $\times1.12$ (4-way), $\times1.14$ (8-way) relative to 1-way.
\end{itemize}
\begin{center}
\small
\begin{tabular}{C{2.4cm}C{2.2cm}C{2.2cm}C{2.2cm}C{2.2cm}}
\toprule
\tblhead{\textbf{Size (KiB)} & \textbf{1-way} & \textbf{2-way} & \textbf{4-way} & \textbf{8-way}}\\
\midrule
4 & 3.44 & 3.25 & \textbf{3.22} & 3.28 \\
8 & 2.69 & 2.58 & \textbf{2.55} & 2.62 \\
16 & \textbf{2.23} & 2.40 & 2.46 & 2.53 \\
32 & \textbf{2.06} & 2.30 & 2.37 & 2.45 \\
64 & \textbf{1.92} & 2.14 & 2.18 & 2.25 \\
128 & \textbf{1.52} & 1.84 & 1.92 & 2.00 \\
256 & \textbf{1.32} & 1.66 & 1.74 & 1.82 \\
512 & \textbf{1.20} & 1.55 & 1.59 & 1.66 \\
\bottomrule
\end{tabular}
\end{center}
AMAT with the clock-stretch penalty included (the bold cell is the best per row; in the slide the \emph{worse-than-left} cells are highlighted). \textbf{Reading:} for tiny (4--8\,KB) caches 2--4 way wins; \textbf{from 16\,KB upward direct-mapped has the lowest AMAT} once the slower clock is charged. This is the quantitative reason an L1 that sets the clock is small and low-associativity, while an L2 (off the critical path) can be highly associative.

\subsubsection{Way prediction}
Keep extra bits per set predicting which way will hit next; the multiplexer is set early and \textbf{only one tag comparison is done} first. A wrong prediction costs extra cycle(s) to check the remaining ways. Alpha 21264 uses it in its 2-way I-cache with $>85\%$ accuracy (SPEC95). Predictor styles: last-used way, MRU, history table. Metrics: accuracy, misprediction rate, energy reduction.
\begin{pick}
Set-associative caches where \textbf{energy and hit time} matter: PMDs, instruction caches (sequential access makes prediction accurate), cores where a full parallel tag compare is too slow/hot. It keeps most of the associativity benefit while activating one way. Cost: predictor hardware, an extra access on a misprediction.
\end{pick}

\subsubsection{Pseudo-associative (column-associative) cache}
Start as a \textbf{direct-mapped} cache; on a miss, probe a second location (invert the high-order index bit) \textbf{sequentially}, not in parallel. Gives fast hit (primary) and slow hit (alternate); swap blocks so the frequent one sits in the fast slot. \textbf{Problem:} a variable 1-or-2-cycle hit is hard for a CPU pipeline $\Rightarrow$ better for \textbf{caches not tied directly to the processor (L2)}; used in the L2 of the MIPS R10000 (and similar in UltraSPARC).
\begin{ex}[AMAT, slide values]
2\,KB: 1-way 5.90, 2-way 4.90, pseudo 4.95. 128\,KB: 1-way 1.50, 2-way 1.45, pseudo 1.371. Pseudo-associativity gets (nearly) 2-way miss rate at direct-mapped hit speed.
\end{ex}


\subsubsection{Compiler optimisations (zero hardware cost)}
\begin{itemize}
\item \textbf{Code:} reorder procedures (profile-guided) to cut conflicts; McFarling: $-50\%$ misses for a 2\,KB direct-mapped I-cache with 4\,B blocks, $-75\%$ for 8\,KB; optimised programs on a direct-mapped cache miss less than unoptimised on 8-way of the same size.
\item \textbf{Merging arrays:} combine parallel arrays (\texttt{val[]}, \texttt{key[]}) into an array of structs $\Rightarrow$ better spatial locality, fewer conflicts.
\item \textbf{Loop interchange:} change loop order so the inner loop walks memory sequentially (row-major). $1024\times1024$ ints, 64\,B lines (16 ints): \textbf{3{,}145{,}728 $\rightarrow$ 196{,}608 misses for $C=A+B$ ($-93.75\%$)}.
\item \textbf{Loop fusion:} merge loops over the same data so a value is reused while in cache (temporal locality). Example: A, B, C larger than the cache: $262{,}144 \rightarrow 196{,}608$ line transfers ($-25\%$).
\item \textbf{Blocking (tiling):} operate on $B\times B$ sub-matrices that fit in cache. Worst-case capacity misses $2N^3+N^2 \rightarrow 2N^3/B+N^2$; also helps register allocation. Slide example: $N=1024$, 32\,KB cache, $B=32$ (three tiles $=3\times4$\,KB$=12$\,KB fits): $\approx1.07$\,billion misses $\rightarrow\approx2.1$ million (simplified model).
\end{itemize}
\begin{pick}
Array/loop-dominated code with large data sets (scientific, image, matrix); embedded systems where adding hardware is impossible. \emph{Matrix multiply with arrays accessed both by rows and columns: interchange does not help $\Rightarrow$ use blocking.}
\end{pick}
\begin{avoid}
Pointer-chasing, irregular data structures (no static pattern); loops where interchange/fusion is illegal because of dependences.
\end{avoid}
\src{ca05 ``Reducing Miss Rate'' through ``Summary of Compiler Optimizations'', ``Way Prediction'', ``Pseudo-Associative Caches''; CachePerformance2 (way prediction, pseudo-associative, loop interchange/fusion/blocking numerics)}

\subsection{Reducing penalty/rate through parallelism}
\subsubsection{Non-blocking (lockup-free) cache}
Cache continues to supply hits during a miss. \textbf{Hit-under-miss} (one outstanding miss), \textbf{hit-under-multiple-miss / miss-under-miss} (overlap several misses; Pentium Pro allows 4). \textbf{Requires an out-of-order CPU} (scoreboard/Tomasulo) and \textbf{multiple memory banks}; greatly increases the cache controller complexity; must detect repeat misses to a block already being fetched.
\begin{itemize}
\item OoO CPUs can hide L1 misses that hit in L2, but \textbf{not a significant fraction of L2 misses}.
\item Example (8\,KB, FP): miss rate$\times$penalty: DM $11.4\%\times16=1.84$; 2-way $10.7\%\times16=1.71$. For integer: DM $7.4\%\times16=1.18$, 2-way $6.0\%\times16=0.96$. Hit-under-miss does \textbf{not increase hit time}, unlike associativity.
\item Hard to evaluate: effective miss penalty is the \emph{non-overlapped} time the CPU stalls.
\end{itemize}
\begin{pick}
Any out-of-order superscalar, especially memory-intensive FP code; pair with banked/interleaved memory. Useless on a strict in-order blocking CPU.
\end{pick}

\subsubsection{Hardware prefetching (stream buffers)}
Fetch what you expect to need before it is requested, outside the cache. Alpha AXP 21064: on an I-miss fetch 2 blocks (target to I-cache, next to the \textbf{instruction stream buffer}); 1/4/16-block ISB catches 15--25\%/50\%/72\% of misses. Data: one data stream buffer removed 25\% of a 4\,KB cache's misses, 4 streams 43\%; 8 streams 50--70\% (for two 64\,KB 4-way caches, scientific programs). \textbf{Needs spare memory bandwidth} (otherwise it hurts).
\begin{ex}[Slide]
Hit time 2\,cc, prefetch-hit time 1\,cc, prefetch-hit rate 25\%, miss rate 1.10\% (8\,KB), miss penalty 50\,cc:\\
$\AMAT=2+0.011\times0.25\times1+0.011\times0.75\times50=2.41525$\,cc. A cache without prefetching needs a miss rate of only 0.83\% (between 8\,KB at 1.10\% and 16\,KB at 0.64\%) to match.
\end{ex}

\subsubsection{Compiler-controlled prefetching}
Special prefetch instructions (register prefetch: HP PA-RISC; cache prefetch: MIPS IV, PowerPC, SPARC v9), non-faulting (a form of speculative execution), best for loops (e.g.\ \texttt{prefetch(b[j+7][0])}). Cost: issuing prefetches takes instruction slots, so benefit must exceed issue cost; needs a non-blocking cache. Works for instructions too.

\begin{dec}[Hardware vs.\ compiler prefetch]
Hardware (stream buffers): sequential/streaming access, no compiler support, spare bandwidth. Compiler: regular strided loops where you know the distance; embedded systems with controllable code. Skip prefetching on bandwidth-starved or energy-limited designs.
\end{dec}
\src{ca05 ``Reducing Cache Miss Penalty or Miss Rate Via Parallelism'' through ``Compiler-Controlled Prefetching''}

\subsection{Reducing hit time}
\begin{itemize}
\item \textbf{Small and simple caches:} the slow part of a hit is index $\rightarrow$ read tag $\rightarrow$ compare. Keep L1 small enough to be on-chip; direct-mapped is simpler/faster (more misses but faster hit); for L2 keep tags on-chip and data off-chip. Used when the cache sets the clock cycle.
\item \textbf{Avoid address translation during indexing.} \emph{Physically addressed} caches translate with the TLB first; a \emph{virtually addressed} cache indexes in parallel with translation. Problems of pure virtual caches:
 \begin{itemize}
 \item \textbf{Protection} bits must be copied on miss and checked every access.
 \item \textbf{Process switch} $\Rightarrow$ same virtual address $\neq$ same physical data $\Rightarrow$ flush (huge task-switch cost and compulsory misses); fix: \textbf{PID in the tag} (PIDs save 0.6--4.3\% miss rate over purging; cost 0.3--0.5\% for uniprocess).
 \item \textbf{Synonyms/aliases:} two virtual addresses for one physical address $\Rightarrow$ two copies; fix: hardware anti-aliasing (check PA on miss, invalidate), or \textbf{page colouring} (aliases share low-order bits, e.g.\ 18 bits).
 \item \textbf{I/O} uses physical addresses; needs mapping to VA.
 \end{itemize}
 Compromise from the slides: overlap cache access with VA translation, requiring the cache index to remain invariant across translation (trivial if the cache is small); used in Alpha 21164 and UltraSPARC III.
\item \textbf{Pipelined cache access / pipelined writes:} write hits normally need tag check before data write; pipeline into two stages so the \emph{current write's tag check overlaps the previous write's data update}; looks like one write per cycle, short cycle time. Works if the CPU does not immediately depend on the written data (watch read/write ordering). Reads are unaffected (tag check in parallel with data read).
\item \textbf{Trace cache:} stores dynamic instruction \emph{traces} from branch entry to exit, avoiding wasted space in conventional blocks (a taken branch every 5--10 instructions means much of a 16-instruction/64\,B block is unused). Costs: complicated address mapping, duplicate storage of the same instructions in different traces. Used in Pentium 4 (NetBurst).
\end{itemize}
\begin{pick}
Small simple L1 / pipelined writes: every high-clock design. Virtual/overlapped translation: tight L1 designs with a small cache. Trace cache: wide-issue front ends that need fetch bandwidth across taken branches (not energy-constrained designs).
\end{pick}
\src{ca05 ``Reducing Hit Time'' through ``Trace Caches (Cont.)''}

\subsection{Cache optimisation summary table (from the slides)}
\small
\begin{longtable}{L{4.6cm}C{1.2cm}C{1.2cm}C{1.2cm}C{1.8cm}L{5.2cm}}
\toprule
\tblhead{\textbf{Technique} & \textbf{Miss pen.} & \textbf{Miss rate} & \textbf{Hit time} & \textbf{HW compl.} & \textbf{Comment}}\\
\midrule
\endhead
Multilevel caches & $+$ & & & 2 & Costly; widely used \\
Critical word first / early restart & $+$ & & & 2 & Widely used \\
Read priority over writes & $+$ & & & 1 & Trivial for uniprocessor \\
Merging write buffer & $+$ & & & 1 & Used with write-through \\
Victim caches & $+$ & $+$ & & 2 & AMD Athlon: 8 entries \\
Larger block size & $-$ & $+$ & & 0 & Trivial; Pentium 4 L2 128\,B \\
Larger cache size & & $+$ & $-$ & 1 & Widely used, esp.\ L2 \\
Higher associativity & & $+$ & $-$ & 1 & Widely used \\
Way-predicting caches & & & $+$ & 2 & UltraSPARC III I-cache \\
Pseudo-associative & & $+$ & & 2 & MIPS R10000 L2 \\
Compiler techniques (miss rate) & & $+$ & & 0 & Software is the challenge \\
Non-blocking caches & $+$ & & & 3 & All out-of-order CPUs \\
HW prefetch (instr./data) & $+$ & $+$ & & 2/3 & UltraSPARC III \\
Compiler-controlled prefetch & $+$ & $+$ & & 3 & Needs non-blocking cache \\
Small and simple caches & & & $+$ ($-$ miss rate) & 0 & Trivial; widely used \\
Avoid translation in indexing & & & $+$ & 2 & Trivial if small; Alpha 21164 \\
Pipelined cache access & & & $+$ & 1 & Widely used \\
Trace cache & & & $+$ & 3 & Pentium 4 \\
\bottomrule
\end{longtable}
\normalsize
($+$ improves, $-$ hurts; HW complexity 0 = none, 3 = highest.) Note the \emph{conflicts of interest}: larger blocks hurt penalty; larger size/associativity hurt hit time; small/simple caches hurt miss rate.

\subsection{Main memory}
\subsubsection{DRAM vs.\ SRAM and key metrics}
\begin{itemize}
\item \textbf{Main memory is DRAM:} dynamic, needs periodic refresh (about 8\,ms, $\sim$1\% of time); address split into RAS and CAS halves (2-D array). \textbf{Caches are SRAM:} no refresh; 6 transistors/bit vs.\ 1. DRAM/SRAM size ratio 4--8; SRAM/DRAM cost and cycle-time ratio 8--16.
\item \textbf{Latency:} access time (request $\rightarrow$ word arrives) and cycle time (minimum between requests, $>$ access time) $\Rightarrow$ miss penalty. \textbf{Bandwidth:} bytes per unit time $\Rightarrow$ time to transfer the block (matters for I/O and large-block L2 misses).
\item DDR trend (slide tables): total best-case access time improved only from 42\,ns (DDR1, 2000) to 26\,ns (DDR4, 2016), whereas DIMM bandwidth rose from 2128 (PC2100) to 21{,}300 MiB/s (PC21300) --- again \textbf{bandwidth far outgrows latency}.
\end{itemize}

\subsubsection{Three ways to raise memory bandwidth}
\begin{center}
\small
\begin{tabular}{L{2.4cm}L{4.6cm}L{4.2cm}L{4.4cm}}
\toprule
\tblhead{\textbf{Organisation} & \textbf{What it is} & \textbf{Pros} & \textbf{Cons}}\\
\midrule
Simple & CPU, cache, bus, memory all one word wide & Cheapest & Block transfer is serial: penalty $\propto$ words per block \\
Wide & Memory/bus (and cache side) $N$ words wide (Alpha 64/256\,bit, UltraSPARC 512\,bit) & Bandwidth $\times N$; penalty drops accordingly & Wider bus cost; MUX between cache and CPU can be on critical path (put MUX between L1 and L2); ECC harder (partial update needs read of the rest); expansion in larger increments \\
Interleaved & Same width bus, $N$ memory banks; bank $=$ address mod $\#$banks & Overlaps bank accesses, shares bus; great for \textbf{sequential} access (cache blocks); cheap & Needs $\ge$ (access time in cc) banks to deliver one word/cycle; many small chips costly; hard to expand \\
Independent banks & Multiple controllers & Multiprocessor/I/O, non-blocking caches (hit under $n$ misses) & Controller complexity \\
\bottomrule
\end{tabular}
\end{center}

\begin{ex}[Slide: what interleaving/width buys]
Base CPI 2, 1.2 mem.\ refs/instr, miss penalty = 4 (address) + 56 (access) + 4 per word. Miss rates: 3\%/2\%/1.2\% for 1/2/4-word blocks.
\begin{center}
\small
\begin{tabular}{L{6.2cm}C{2.3cm}C{2.3cm}C{2.3cm}}
\toprule
\tblhead{\textbf{Configuration} & \textbf{1-word} & \textbf{2-word} & \textbf{4-word}}\\
\midrule
64-bit bus \& memory, no interleaving & 4.30 & 5.07 & 5.69 \\
64-bit bus \& memory, interleaved & --- & 3.63 & \textbf{3.09} \\
128-bit bus \& memory, no interleaving & --- & \textbf{3.54} & 3.84 \\
\bottomrule
\end{tabular}
\end{center}
\textit{Note:} the slide labels the narrow bus ``64-bit'' and the wide bus ``128-bit''; its arithmetic treats narrow as one word per transfer and wide as two words per transfer. \textbf{Reading:} wide/interleaved beat simple by a large margin; with larger blocks \textbf{interleaving} wins (it overlaps accesses on a narrow bus), with 2-word blocks a \textbf{128-bit} bus is marginally better. If the bus cannot be widened (pins/cost), choose interleaving.
\end{ex}
\src{ca05 ``Main Memory'' through ``Independent Memory Banks''}

\subsubsection{Memory hierarchy by device class (slide figures)}
\textbf{PMD:} L1 64\,KB (1\,ns), L2 256\,KB (5--10\,ns), no L3, 1--2\,GB DRAM (50--100\,ns), flash 4--64\,GB. \textbf{Laptop/desktop:} adds L3 (4--8 / 8--32\,MB, 10--20\,ns), 4--64\,GB DRAM, flash. \textbf{Server:} L3 16--64\,MB, 32--256\,GB DRAM, flash 1--16\,TB and disk 16--64\,TB (5--10\,ms).

\subsection{Virtual memory}
\begin{itemize}
\item \textbf{Why:} many processes each with their own address space share a smaller physical memory; need \emph{protection} (a process only touches its own blocks); faster program start (not everything resident); frees the programmer from overlays.
\item \textbf{Pages vs.\ segments:} pages $=$ fixed-size blocks (4--8\,KB), one address (page number $+$ offset), \textbf{no external fragmentation}, allocation anywhere (non-contiguous), no compaction; segments $=$ variable size (1\,B up to $2^{16}$--$2^{32}$\,B), address needs 2 words (segment number + offset). \textbf{Paging is the default choice}.
\item \textbf{Translation:} virtual page number indexes the \textbf{page table} (valid bit, physical page number); valid $=0\Rightarrow$ \textbf{page fault} (OS loads from disk). The page table is in memory, so every access would cost two memory accesses $\Rightarrow$ \textbf{TLB}, a special cache of translations.
\item \textbf{TLB:} usually small ($\le128$--256 entries); high-end: fully associative; mid-range: small $n$-way set associative. Fields: valid, dirty, tag, physical page number. Physical address then drives cache tag compare (TLB $\rightarrow$ cache interaction).
\item Example split: 32-bit VA $=$ 20-bit VPN $+$ 12-bit offset (4\,KB pages).
\end{itemize}
\src{ca05 ``Virtual Memory'' through ``Techniques for Fast Address Translation''}

\subsection{Scenario $\rightarrow$ memory design, quick reference}
\small
\begin{longtable}{L{5.2cm}L{4.1cm}L{6.3cm}}
\toprule
\tblhead{\textbf{Symptom / constraint} & \textbf{Choose} & \textbf{Reason}}\\
\midrule
\endhead
Small direct-mapped L1, ping-pong conflicts & Victim cache; or pseudo-assoc./2-way; or array merging/padding & Conflict misses; victim cache keeps DM hit time. \\
Streaming/sequential large arrays & Stream-buffer prefetch, moderate block (64\,B), interleaved memory, loop interchange & Compulsory/spatial misses; bandwidth is plentiful. \\
Large matrix multiply & Blocking (tiling) with 3 tiles $\le$ cache & Capacity misses; interchange cannot help both row/column accesses. \\
OoO core with many outstanding loads & Non-blocking cache + banked memory & Overlap misses; otherwise OoO core is wasted. \\
L1 hit time limits clock & Small, simple (DM) L1, pipelined writes, VA/translation overlap & Hit time is multiplied across every instruction. \\
Write-through L1 & Merging write buffer + read priority over writes & Prevent buffer-full stalls and RAW delays. \\
Large blocks and long block fill & Critical word first \& early restart & CPU resumes when the needed word arrives. \\
Instruction-fetch bound, taken branches every 5--10 instr. & Trace cache or I-stream buffer; BTB & Fetch bandwidth across taken branches. \\
Cache must be tiny but fast & Direct-mapped, small block & Table: DM wins AMAT $\ge$16\,KB; for 4--8\,KB use 2--4 way. \\
\bottomrule
\end{longtable}
\normalsize



\section{Thread-Level Parallelism: Multithreading, SMT, Multicore}\label{ch:tlp}

\subsection{Decision map}
\small
\begin{tabularx}{\textwidth}{L{4.4cm}L{4.2cm}X}
\toprule
\tblhead{\textbf{Scenario} & \textbf{Feature} & \textbf{Why}}\\
\midrule
Throughput server, many independent threads/requests & Multicore (CMP) with multithreading per core & Applications with TLP (DB, web) map directly; ILP alone is saturated. \\
Pipeline stalls are \emph{long and rare} (e.g.\ L2 misses), want no slowdown of a single thread & Coarse-grained multithreading & Switch only on costly stalls; pipeline refill $\ll$ stall time. \\
Stalls are \emph{short and frequent}; hardware can switch every cycle & Fine-grained multithreading & Hides short and long stalls; interleaves each cycle. \\
Out-of-order superscalar with idle issue slots; want cheap extra throughput & SMT & Reuses renaming and OoO hardware; fills unused slots from another thread. \\
Single important thread must run as fast as possible & One big (SMT-capable) superscalar core; preferred-thread priority & Multi-core does not speed one thread; SMT with priority protects it. \\
Limited die area / thermal budget but threaded software & Several small cores & Each core simpler and easier to design/manufacture; better energy per throughput. \\
Applications hard to parallelise & Do \emph{not} rely on cores/threads & TLP gives no gain; spend on ILP/caches. \\
\bottomrule
\end{tabularx}
\normalsize

\subsection{Why TLP at all}
\begin{itemize}
\item ILP gives implicit parallelism within a loop or basic block; it saturates at 3--6 issue (Section~\ref{ch:ilp}), and for most applications \textbf{most execution units sit idle} even on an 8-way superscalar (Tullsen et al.).
\item \textbf{TLP is explicit:} multiple threads, each with its own instructions, data, PC and register state. Goals: improve \emph{throughput} of machines running many programs and \emph{execution time} of multithreaded programs. TLP ``could be more cost-effective to exploit than ILP''.
\item Natural parallelism: online transaction systems (many independent queries/updates), data mining, databases, scientific codes; server-per-client threads; games (AI, graphics, physics as threads).
\item Industry move: IBM Power4 (2000) was the first commercial single-chip general-purpose multiprocessor (2 Power3 + shared L2); Sun, AMD, Intel followed.
\end{itemize}
\src{ca06 ``Performance beyond single thread ILP'', ``Thread Level Parallelism (TLP)''; ca08 ``Why multi-core?'', ``Thread-level parallelism''}

\subsection{Multithreading on one core}
Multiple threads \textbf{share the functional units} of one processor by overlapping. The processor \emph{duplicates the independent state of each thread}: a separate register file, a separate PC, and (for independent programs) a separate page table. Memory is shared via the virtual-memory mechanism. Hardware thread switch is \emph{much faster} than a full process switch (100s--1000s of clocks).

\begin{center}
\small
\begin{tabular}{L{2.4cm}L{4.2cm}L{4.5cm}L{4.4cm}}
\toprule
\tblhead{\textbf{} & \textbf{Fine-grained} & \textbf{Coarse-grained} & \textbf{SMT (simultaneous)}}\\
\midrule
When it switches & Every instruction/clock, round-robin, skipping stalled threads & Only on costly stalls (e.g.\ L2 miss) & No switching: instructions from several threads issue in the \emph{same cycle} \\
Hides & Short \emph{and} long stalls & Only long stalls (pipeline refill $\ll$ stall time) & Short and long stalls \emph{and} wasted issue slots \\
Single-thread cost & Slows a ready thread (delayed by others) & Does not slow a thread (others run only when it stalls) & Can slow a thread unless priority is given \\
Hardware & Switch every cycle; per-thread state & No very fast switch needed; pipeline must be flushed/frozen on switch & OoO superscalar + per-thread rename table, PC, (logically) per-thread ROB \\
Example & Sun Niagara & IBM AS/400 & Pentium 4 (HT), IBM Power5 \\
\bottomrule
\end{tabular}
\end{center}

\begin{dec}[Pick the multithreading style]
\begin{itemize}
\item \textbf{Coarse-grained} if the core is simple/in-order, stalls are dominated by rare, long events (L2/L3 misses), and single-thread speed must not degrade; avoid when short stalls dominate (refill cost cannot be recovered).
\item \textbf{Fine-grained} for throughput-oriented cores with many threads where single-thread latency is secondary (server/web/DB); needs per-cycle switching hardware.
\item \textbf{SMT} when you already have a wide out-of-order superscalar (renaming, large physical register set, OoO completion): \emph{the marginal hardware is small}. Not worthwhile for narrow in-order cores.
\end{itemize}
\end{dec}

\subsubsection{Why SMT is cheap on an OoO core}
Dynamically scheduled processors already have: a large set of virtual (physical) registers that can hold several threads' register sets; \textbf{register renaming} gives unique register ids so instructions from different threads mix in the datapath without confusing sources and destinations; \textbf{out-of-order completion} lets threads proceed independently. You only add a per-thread renaming table and separate PCs; independent commit is obtained by logically keeping a separate ROB per thread.

\subsubsection{SMT design challenges and costs (IBM Power5 case study)}
\begin{itemize}
\item \textbf{Single-thread performance:} SMT needs fine-grained scheduling; a \textbf{preferred-thread} approach protects one thread but loses some throughput when that thread stalls. Power5 gives hardware-controllable \emph{relative thread priority}; with balanced threads both run slower than if each owned the machine.
\item Larger register file for multiple contexts; do not lengthen the clock cycle (more candidate instructions to issue, harder choice at commit); avoid cache/TLB conflicts between threads.
\item \textbf{Why only 2 threads in Power5?} With 4, one of the shared resources (physical registers, cache, memory bandwidth) would become a bottleneck.
\item \textbf{Changes made in Power5 to support SMT:} higher associativity in the L1 I-cache and instruction translation buffers; per-thread load/store queues; larger L2 (1.92 vs.\ 1.44\,MB) and L3; separate instruction prefetch and buffering per thread; physical registers 152 $\rightarrow$ 240; larger issue queues. \textbf{Power5 core is $\approx$24\% larger than Power4's} because of SMT.
\end{itemize}

\subsubsection{SMT performance data to quote}
\begin{center}
\small
\begin{tabular}{L{6.2cm}L{9.2cm}}
\toprule
\tblhead{\textbf{Measurement} & \textbf{Result (slides)}}\\
\midrule
Pentium 4 Extreme SMT, SPECrate (same benchmark $\times$ copies) & 1.01 (SPECint\_rate), 1.07 (SPECfp\_rate) \\
Pentium 4, 26 SPEC benchmarks paired with each other (262 runs) & Speedups 0.90--1.58, \textbf{average 1.20} \\
Power5, 8-processor server & 1.23 (SPECint\_rate), 1.16 (SPECfp\_rate) with SMT \\
Power5, two copies of each application & 0.89--1.41; most gained; \textbf{FP applications had the most cache conflicts and the least gains} \\
General statement (multicore deck) & ``Enables better threading (e.g.\ up to 30\%)'' \\
\bottomrule
\end{tabular}
\end{center}
\textbf{Reading:} SMT gives a modest, workload-dependent gain (roughly 1.0--1.4$\times$) and can \emph{hurt} (speedup $<$1) when threads fight over caches. It is \emph{not} a true parallel processor: there is a single copy of each resource, and threads cannot use the same functional unit simultaneously.

\begin{ex}[Area-efficiency argument for SMT vs.\ a second core]\bs
Power5 SMT cost $\approx$24\% extra core area for $\approx$16--23\% throughput. A second core costs about 100\% extra core area but can give up to $\sim$2$\times$ throughput for fully parallel software. So SMT is roughly \emph{area-neutral} on throughput but needs little extra design and shines when threads are stall-bound; a second core wins when software is truly parallel and area/power allow.
\end{ex}
\src{ca06 ``New Approach: Multithreaded Execution'' through ``Initial Performance of SMT''; ca08 ``Simultaneous multithreading (SMT)''}

\subsection{Multicore (chip multiprocessor, CMP)}
\begin{itemize}
\item Several cores on one processor socket; they run in parallel; \textbf{within each core threads are time-sliced} as on a uniprocessor. The OS sees each core as a separate processor and schedules threads/processes across them.
\item \textbf{Why multicore:} single-core clock frequencies cannot rise further; limited ILP; many applications are already multithreaded; general trend toward parallelism.
\item A multicore is a \textbf{shared-memory MIMD multiprocessor on one chip}: different cores run different instruction streams on different data, all sharing the same memory.
\item \textbf{Applications that benefit:} database servers, web servers, compilers, multimedia, scientific, CAD/CAM; ``anything that can be threaded''. Some applications are difficult to parallelise.
\end{itemize}

\subsubsection{Multicore vs.\ SMT, and combining them}
\begin{center}
\small
\begin{tabular}{L{3.2cm}L{6.0cm}L{6.2cm}}
\toprule
\tblhead{\textbf{} & \textbf{Multicore} & \textbf{SMT}}\\
\midrule
Resources & Each core has its own copy of resources & Single copy of each resource shared by threads \\
Core size/design & Several smaller, less powerful cores; easier to design and manufacture & One large, fast superscalar core \\
Strength & Great with thread-level parallelism; independent execution & Great single-thread performance; still mostly ILP-based \\
Weakness & Does nothing for one thread; coherence traffic & Not true parallelism; shared-resource contention; $\le$30\% gain \\
Threads & One (or more) per core & 2, 4 or sometimes 8 per core (Intel ``hyper-threads'') \\
\bottomrule
\end{tabular}
\end{center}
Four combinations: single-core non-SMT (standard uniprocessor), single-core SMT, multi-core non-SMT, \textbf{multi-core with SMT} (the best-throughput option; e.g.\ SMT dual-core runs four threads concurrently).

\subsubsection{Memory hierarchy in threaded designs}
\begin{itemize}
\item SMT only $\Rightarrow$ all caches shared. Multicore $\Rightarrow$ \textbf{L1 caches private}; L2 \textbf{private in some designs, shared in others}; \textbf{main memory always shared}. Example: dual-core hyper-threaded Xeon with private L1 and \textbf{shared} L2.
\item \textbf{Shared L2:} threads share capacity dynamically and share data without duplication\bs; risk of interference. \textbf{Private L2:} isolation and simple timing\bs; duplicated data, wasted capacity if load is unbalanced\bs. In either case sharing introduces \textbf{coherence misses} (true/false sharing; Section on miss types).
\end{itemize}

\begin{pick}[Multicore + SMT]
Server/throughput workloads and mixed multitasking desktops: plenty of independent threads, area available, TLP-friendly software. Add SMT to each core if the core is a wide OoO and memory stalls are common.
\end{pick}
\begin{avoid}
Single-threaded or hard-to-parallelise applications (give one fast core instead); extremely tight power budgets that cannot afford duplicate structures or shared-cache contention.
\end{avoid}

\begin{dec}[ILP vs.\ TLP balance (slide conclusion)]
There is no universal answer: \emph{server} market $\Rightarrow$ more TLP (many cores/threads); \emph{desktop} where single-thread speed is still a primary requirement $\Rightarrow$ keep a strong ILP core but add a few cores. ``Balance of ILP and TLP decided in marketplace.''
\end{dec}

\begin{trap}[TLP traps]
\begin{itemize}
\item SMT threads cannot use the \emph{same} functional unit at the same time; it is not a doubling of resources.
\item Coarse-grained multithreading is useless when stalls are short (pipeline refill cost dominates); fine-grained slows individual threads.
\item More threads $\ne$ more speed: Power5 stops at 2 threads because shared resources saturate; SMT can give speedup $<1$ on cache-conflict-heavy (FP) code.
\item Multicore does not accelerate a single thread; the OS must have threads to schedule.
\item False sharing causes coherence misses even though data is not shared.
\end{itemize}
\end{trap}
\src{ca08 (multicore deck: ``Multi-core architectures'' through ``Designs with private L2 caches''); ca06 ``Summary'', ``Conclusions''; Cache-Misses (coherence misses)}


\section{Data-Level Parallelism: Vector, SIMD and GPU (SIMT)}\label{ch:dlp}
\textit{Source: Hennessy \& Patterson, Ch.~4 (not in the provided slide decks, but the professor named SIMT/DLP as examinable). Standard textbook material; no daggers used inside this section.}

\subsection{Decision map}
\small
\begin{tabularx}{\textwidth}{L{4.4cm}L{4.3cm}X}
\toprule
\tblhead{\textbf{Scenario} & \textbf{Feature} & \textbf{Why}}\\
\midrule
Same operation on large regular arrays (image, matrix, signal, ML) & DLP hardware: vector unit, SIMD extension or GPU & One instruction fetch/decode controls many operations: far less control energy and area per operation than MIMD cores. \\
Long vectors, scientific, need predictable bandwidth & Vector architecture (vector registers, lanes, chaining, many memory banks) & Deep pipelined lanes, strided/gather loads in hardware, compact code. \\
General-purpose CPU, short vectors (media, crypto, small kernels) & SIMD extension (SSE/AVX/NEON) & Cheap to add to an existing core; no data transfer to another device. \\
Massive parallelism (10$^5$+ independent elements), throughput over latency & GPU (SIMT, many SMs, many warps) & Thousands of lanes; memory latency hidden by switching warps instead of big caches/OoO. \\
Branchy per-element code, pointer structures, little data & Multicore CPU (MIMD) & Each core has its own PC: no divergence penalty; caches and OoO handle irregularity. \\
Data must be shipped to an accelerator for little work & Keep it on the CPU (SIMD) or batch the work & Transfer time can exceed compute time. \\
\bottomrule
\end{tabularx}
\normalsize

\subsection{Why DLP is efficient, and how to spot it in the workload}
\begin{itemize}
\item \textbf{Efficiency:} one instruction describes $N$ operations, so fetch, decode, dependence checking and scheduling are paid once per $N$ operations; elements of one vector instruction are independent by definition (no hazard checks between them); the memory access pattern is known, so it can be pipelined and banked.
\item \textbf{Spotting DLP:} a loop whose iterations are \emph{independent} (\texttt{y[i] = a*x[i] + y[i]}). A \textbf{loop-carried dependence} (\texttt{a[i] = a[i-1] + b[i]}) blocks it, except two special cases: a \textbf{reduction} (\texttt{sum += a[i]}) is parallelised as a tree in $\log_2 N$ steps, and a \textbf{scan/prefix-sum} has parallel algorithms. Conditional code inside the loop becomes \emph{masked} execution.
\item \textbf{Write in the answer:} ``the workload has DLP because iteration $i$ only touches element $i$; the order of execution between elements does not matter; the only ordering is the final reduction''.
\end{itemize}

\subsection{Vector architectures}
\begin{itemize}
\item \textbf{Vector registers} hold many elements (e.g.\ VMIPS: 8 registers $\times$ 64 elements $\times$ 64 bits). \textbf{Vector length register (VL)} handles $n<$ maximum vector length (MVL); \textbf{strip mining} handles $n>$ MVL by looping over MVL-sized chunks.
\item \textbf{Lanes:} the functional unit is replicated $L$ times; element $i$ goes to lane $i \bmod L$. Doubling lanes halves time per vector instruction without changing the code.
\item \textbf{Chaining:} a dependent vector instruction starts as soon as the first result element is ready (vector version of forwarding).
\item \textbf{Timing model:} group instructions that can run together into \emph{convoys}; one convoy takes one \emph{chime} ($\approx n/L$ cycles). Time $\approx$ (number of convoys) $\times\, n/L$ + start-up latency.
\item \textbf{Vector mask registers} implement \texttt{if} inside loops (operate only where the mask bit is 1; wasted lanes where it is 0).
\item \textbf{Memory:} must deliver one element per lane per cycle, so use \textbf{many independent banks}. Unit stride: need $\#\text{banks} \ge$ bank busy time. Non-unit stride $s$ conflicts when $\#\text{banks}/\gcd(s,\#\text{banks}) <$ bank busy time. Sparse data uses \textbf{gather/scatter} (index-vector loads/stores), which are much slower than unit stride.
\end{itemize}

\subsection{SIMD multimedia extensions (SSE, AVX, NEON)}
Fixed-width registers (128/256/512 bits) split into packed elements (e.g.\ 8 floats in 256 bits). \textbf{Pros:} small hardware cost, uses the existing caches and core, no transfer. \textbf{Cons:} width fixed in the instruction set (new instructions per generation), historically weak masking and gather (improved in AVX-512), alignment matters, compilers vectorise only simple loops. Peak speedup $\le$ number of elements per register.

\subsection{GPUs and SIMT (Single Instruction, Multiple Threads)}\label{sec:simt}
\textbf{Programming model:} the programmer writes code for \emph{one} scalar thread; thousands of threads are launched as a \textbf{grid} of \textbf{thread blocks}. \textbf{Hardware:} each block runs on one \textbf{streaming multiprocessor (SM)}; the SM groups 32 threads into a \textbf{warp} that executes one instruction at a time across 32 SIMD lanes in lock-step. Each thread has its own registers and can take its own branch; the hardware handles this with masks.
\begin{itemize}
\item \textbf{Latency hiding by multithreading, not by caches/OoO:} an SM keeps many warps resident (registers of all of them stay in a very large register file). When a warp waits for memory (hundreds of cycles), the \textbf{warp scheduler} issues from another ready warp in the next cycle at zero switch cost. This is \emph{fine-grained multithreading} (\S\ref{ch:tlp}) at large scale. It needs many warps in flight.
\item \textbf{Occupancy} $=$ resident warps / maximum warps per SM. Limited by registers per thread, shared memory per block, and threads per block. Low occupancy $\Rightarrow$ not enough warps to hide latency $\Rightarrow$ SM idles.
\item \textbf{Memory coalescing:} if the 32 threads of a warp access consecutive addresses, the hardware merges them into a few wide transactions. Strided or random addresses become up to 32 separate transactions $\Rightarrow$ bandwidth wasted. Fix: thread $i$ accesses element $i$; struct-of-arrays layout; stage awkward accesses through shared memory.
\item \textbf{Shared memory:} a software-managed scratchpad per SM (static on-chip placement). Use it to load a tile once and reuse it many times (stencils, matrix multiply), raising arithmetic intensity.
\item \textbf{Synchronisation:} barrier \emph{within a block} (\texttt{\_\_syncthreads}); \emph{no} cheap barrier across blocks (end the kernel, or use atomics). Global atomics on one address serialise: privatise in shared memory, then merge (histogram pattern).
\item \textbf{Host--device transfer:} discrete GPUs receive data over PCIe, which is far slower than GPU memory. If compute per byte transferred is small, the transfer dominates: keep data resident on the GPU across kernels, batch work, and overlap transfer with compute (double buffering / streams).
\end{itemize}

\subsubsection*{Branch divergence (the main SIMT inefficiency)}\label{sec:divergence}
If threads of one warp take different paths, the warp executes \emph{both} paths one after the other, masking off the lanes not on that path, and reconverges after the branch.
\[
\text{SIMD efficiency}=\frac{\text{average active lanes}}{32},\qquad
T_{warp}\approx T_{if}+T_{else}\ \text{ (if any lane takes each side)}
\]
\begin{ex}[Divergence in numbers]
30\% of pixels (randomly placed) need a 50-instruction path, the rest a 5-instruction path. Probability a warp has \emph{no} expensive pixel: $0.7^{32}\approx 10^{-5}$, so essentially every warp pays $50+5=55$ instructions, against an average useful work of $0.3\times50+0.7\times5=18.5$. Efficiency $\approx 34\%$. \textbf{Fixes:} (1) \emph{compaction}: a first pass writes the indices of expensive pixels into a list; a second kernel processes only that list with full warps (dynamic, data-dependent decision); (2) sort/bin data so a warp's threads branch the same way; (3) predication for short branches (static, compiler).
\end{ex}

\subsection{Comparison table}
\small
\begin{longtable}{L{2.3cm}L{3.2cm}L{3.2cm}L{3.4cm}L{3.4cm}}
\toprule
\tblhead{\textbf{} & \textbf{Vector} & \textbf{SIMD extension} & \textbf{GPU (SIMT)} & \textbf{Multicore CPU (MIMD)}}\\
\midrule
\endhead
Control & One instr, VL elements & One instr, fixed width & One instr per warp of 32 threads & One PC per core \\
Hides memory latency by & Deep pipelining, banks & Caches, OoO core & Switching among many warps & Caches, OoO, SMT \\
Branches in data loop & Mask registers & Masks/blend & Masks + reconvergence (divergence cost) & Free (independent PCs) \\
Irregular memory & Gather/scatter & Limited gather & Uncoalesced = slow & Caches tolerate it \\
Best for & Long regular vectors & Short vectors on CPU & Huge, regular, throughput & Irregular, branchy, task-parallel \\
Weak at & Short vectors & Width/ISA churn & Transfer, divergence, small problems & Ops per watt on regular data \\
\bottomrule
\end{longtable}
\normalsize

\begin{trap}[DLP traps]
\begin{itemize}
\item A GPU is often fast because of its \emph{memory bandwidth}, not its FLOPs: SAXPY ($y=ax+y$) has AI $= 2$\,FLOP$/12$\,B and is memory-bound on every device (\S\ref{sec:roofline}).
\item More lanes do not fix divergence or uncoalesced access; the waste grows with the width.
\item Loop-carried dependences cannot be vectorised (except reductions/scans, which need a tree and a final sync).
\item GPU caches are small per thread; GPUs rely on parallel slack (many warps), so a GPU with too little work runs badly.
\item Offloading small problems: transfer + launch overhead exceed the compute.
\end{itemize}
\end{trap}

\section{Multiple Processing Elements: Distribution, Timing, Communication, Synchronisation}\label{ch:pe}
\textit{Source: Hennessy \& Patterson Ch.~5 and standard parallel-computing material (the slides cover only multicore basics and coherence misses; the professor named these topics). Lab 5 is direct evidence for this section.}

\subsection{Decision map}
\small
\begin{tabularx}{\textwidth}{L{4.6cm}L{4.4cm}X}
\toprule
\tblhead{\textbf{Scenario} & \textbf{Choice} & \textbf{Why}}\\
\midrule
Independent equal-cost items & Static block partition & Zero scheduling overhead, contiguous data per PE (locality, no sharing). \\
Unequal or unknown task costs & Dynamic work queue / work stealing & Idle PE takes more work; balances automatically. \\
Cost pattern known (e.g.\ heavy tasks clustered at one end) & Static cyclic or block-cyclic & Spreads heavy tasks evenly at no run-time cost. \\
All PEs update a shared result & Private partial results + tree reduction & Removes per-update coherence traffic and locking. \\
Stages producer$\rightarrow$consumer & Pipeline across PEs with queues & Overlaps stages; throughput set by slowest stage, so balance stages. \\
Few cores, shared bus & Snooping coherence (MSI/MESI) & Simple; every cache sees every miss. \\
Many cores / multi-socket & Directory coherence, NUMA-aware placement & Broadcast does not scale; point-to-point messages do. \\
Grid/stencil computation & 2D block partition (not strips) & Less boundary per PE $\Rightarrow$ less communication. \\
Heterogeneous PEs (big + small, CPU + GPU) & Split proportional to speed, or dynamic & Equal split leaves fast PEs idle. \\
\bottomrule
\end{tabularx}
\normalsize

\subsection{Order of execution: read it from the workload}\label{sec:order}
Draw the \textbf{dependence (task) graph}: nodes are tasks/instructions, edges are ``must happen before'' (true data dependence, RAW). Then:
\begin{itemize}
\item \textbf{Work} $T_1$ = total time on one PE; \textbf{span} $T_\infty$ = longest path (critical path). \textbf{Available parallelism} $=T_1/T_\infty$. With $p$ PEs: $T_p \ge \max(T_1/p,\ T_\infty)$. More PEs than $T_1/T_\infty$ are wasted.
\item \textbf{Amdahl} (fixed problem): $S(p)=\dfrac{1}{(1-f)+f/p}$, limit $1/(1-f)$. \textbf{Gustafson} (problem grows with $p$): $S(p)=(1-f)+f\,p$. \textbf{Efficiency} $E=S/p$.
\item Within one core the same idea sets ILP: the recurrence (loop-carried chain) bounds IPC regardless of width (Lab 10: 0.27 IPC vs.\ $\approx$3 for independent chains).
\end{itemize}
\small
\begin{longtable}{L{3.3cm}L{3.0cm}L{4.4cm}L{4.6cm}}
\toprule
\tblhead{\textbf{Pattern in workload} & \textbf{Parallel?} & \textbf{How to execute} & \textbf{Sync / bottleneck}}\\
\midrule
\endhead
Independent iterations / requests & Fully (DLP/TLP) & SIMD/GPU or block partition over PEs & Only at the end; usually memory bandwidth \\
Reduction (sum, max, histogram) & Yes, as a tree & Private partials, combine in $\log_2 p$ steps & One combine phase; never a shared accumulator \\
Loop-carried recurrence ($x_i=f(x_{i-1})$) & No (serial) & One PE; ILP only; restructure the algorithm if possible & Critical path $=n\times$ latency \\
Producer $\rightarrow$ consumer stages & Pipeline parallel & One stage per PE, bounded queues & Throughput $=1/$slowest stage; queue sync \\
Phases separated by global steps & Fork--join & Parallel phase, barrier, next phase & Barrier cost + imbalance per phase \\
Irregular task DAG & Partly & Dynamic scheduler: run any ready task & Scheduler overhead, critical path \\
Shared read-only data & Yes & Replicate in caches & None (Shared state) \\
Shared read-write data & Limited & Partition ownership or lock & Coherence + lock contention \\
\bottomrule
\end{longtable}
\normalsize

\subsection{Distributing the work and its timing}\label{sec:distrib}
\[
T_{par}=\underbrace{T_{serial}}_{\text{Amdahl}}+\max_{i}\Big(\frac{W_i}{\text{speed}_i}\Big)+T_{comm}+T_{sync},\qquad
\text{load imbalance}=\frac{\max_i W_i}{\text{avg}_i W_i}
\]
\begin{itemize}
\item \textbf{Decomposition:} \emph{data} (split the data, every PE runs the same code: SPMD, matches DLP), \emph{task/functional} (different PEs do different jobs), \emph{pipeline} (stages).
\item \textbf{Mapping:} \emph{block} (contiguous chunk per PE: best locality), \emph{cyclic} (item $i\rightarrow$ PE $i \bmod p$: best balance when cost varies smoothly), \emph{block-cyclic} (compromise).
\item \textbf{Static vs.\ dynamic:} static is free at run time and keeps locality but needs predictable costs; dynamic (shared queue, work stealing, guided chunks) adapts but costs a queue access per chunk and may lose locality.
\item \textbf{Granularity:} fine tasks balance well but pay more overhead per unit of work; coarse tasks have low overhead but risk imbalance. Pick chunks so overhead is a few \% of chunk time.
\end{itemize}

\begin{ex}[Static vs.\ dynamic timeline: write this kind of calculation]
4 PEs, 16 tasks: tasks 0--3 cost 4 units, tasks 4--15 cost 1 unit. Total $=28$, ideal $=7$ per PE.
\begin{center}\small
\begin{tabular}{L{4.4cm}L{5.4cm}C{1.3cm}C{1.6cm}}
\toprule
\tblhead{\textbf{Policy} & \textbf{PE loads} & \textbf{$T_{par}$} & \textbf{Speedup}}\\
\midrule
Static block (4 consecutive each) & PE0: 4+4+4+4$=$16; PE1--3: 4 each & 16 & 1.75 \\
Static cyclic ($i \bmod 4$) & each PE: 4+1+1+1$=$7 & 7 & 4.0 \\
Dynamic queue (0.1 overhead/fetch) & heavy tasks first, then 3 unit tasks each & $\approx$7.4 & 3.8 \\
\bottomrule
\end{tabular}
\end{center}
\textbf{Reading:} cyclic is best \emph{only because} we knew where the heavy tasks are (static knowledge). If costs are unknown or input-dependent, dynamic gives near-ideal time without that knowledge, at a small overhead. Write this reasoning in the exam.
\end{ex}

\subsection{Communication}\label{sec:comm}
\begin{itemize}
\item \textbf{Shared memory} (multicore): communicate by loads/stores; hardware coherence moves the data; cheap to program, hidden traffic. \textbf{Message passing} (clusters, some accelerators): explicit send/receive; cost visible, scales further.
\item \textbf{Cost model:} $T_{msg}=$ latency $+$ size/bandwidth. Many small messages are latency-bound: batch them.
\item \textbf{Communication-to-computation ratio} decides if parallelism pays. For an $N\times N$ grid on $p$ PEs: strips exchange $\approx 2N$ boundary values per PE; square blocks exchange $\approx 4N/\sqrt{p}$ (\emph{surface-to-volume}). Prefer blocks for large $p$.
\item \textbf{Overlap} communication with computation (prefetch the next boundary while computing the interior; double buffering).
\item \textbf{NUMA:} memory is attached to sockets; remote access is slower. Allocate data on the node of the PE that uses it (owner computes).
\end{itemize}

\subsection{Cache coherence}\label{sec:coh}
\textbf{Problem:} private caches may hold different values of the same address. \textbf{Coherence} guarantees (1) a read returns the latest write, and (2) all PEs see writes to one location in the same order. Standard approach: \textbf{write-invalidate} (a writer gets the only copy; other copies are invalidated).
\small
\begin{center}
\begin{tabular}{L{2.3cm}L{13.2cm}}
\toprule
\tblhead{\textbf{State} & \textbf{Meaning (MSI; MESI adds E)}}\\
\midrule
M (Modified) & Only copy, dirty; this cache must supply the data and write it back \\
E (Exclusive) & Only copy, clean; can be written \emph{silently} (no bus transaction) $\rightarrow$ M. Saves traffic for private data \\
S (Shared) & Clean, possibly in several caches; a write must first invalidate the others \\
I (Invalid) & Not usable; a read or write causes a miss \\
\bottomrule
\end{tabular}
\end{center}
\normalsize
\begin{itemize}
\item \textbf{Snooping:} every cache watches a shared bus/broadcast. Simple and fast for a few cores; the broadcast becomes the bottleneck as cores grow.
\item \textbf{Directory:} each memory block has an entry listing its sharers (at its home node); requests go point-to-point to the home and only to sharers. Scales to many cores and sockets; costs directory storage and an extra hop.
\item \textbf{Coherence misses} (4th C): \textbf{true sharing} (PEs really exchange the data) and \textbf{false sharing} (different variables in the same cache line, so the line ping-pongs although no data is shared). Fix false sharing by padding/aligning per-PE data to the line size; fix true sharing by privatisation + reduction.
\item \textbf{Larger blocks} increase false sharing; more cores increase invalidation traffic.
\end{itemize}

\subsection{Synchronisation}\label{sec:sync}
\begin{itemize}
\item \textbf{Hardware primitives:} atomic read-modify-write: test-and-set, fetch-and-add, compare-and-swap (CAS), load-linked/store-conditional (LL/SC). Locks, barriers and lock-free updates are built from them.
\item \textbf{Spin lock cost:} spinning with test-and-set writes the lock line repeatedly $\Rightarrow$ coherence storm. \textbf{Test-and-test-and-set} spins on a cached (Shared) copy and attempts the atomic only when the lock looks free. Under contention, $n$ waiting PEs still cause bursts of invalidations at every release; queue locks\bs{} or less sharing fix it.
\item \textbf{A critical section is a serial fraction:} if each of $p$ PEs holds a lock for time $c$ per item, the locked part runs one PE at a time (Amdahl). Shorten it, split the lock (finer granularity: one lock per bucket), or remove it (privatise).
\item \textbf{Barriers} make every PE wait for the slowest: cost $=$ barrier latency $+$ imbalance, paid at every barrier. Use fewer, larger phases.
\item \textbf{Reductions:} each PE accumulates privately; combine in a tree ($\log_2 p$ steps) at the end. One sync instead of one per update.
\item \textbf{Memory consistency:} \emph{sequential consistency} (all PEs see memory operations in one interleaving that respects each program's order) is simple but limits hardware reordering (store buffers). Real machines use \emph{relaxed} models and provide \textbf{fences}/acquire--release; correctly synchronised (data-race-free) programs behave as if sequentially consistent.
\end{itemize}

\begin{center}
\small
\begin{tabular}{L{4.4cm}L{5.4cm}L{5.6cm}}
\toprule
\tblhead{\textbf{Sync problem} & \textbf{Technique} & \textbf{Why better / cost}}\\
\midrule
Global counter updated by all & Per-PE counters (padded) + final reduction & $O(\text{updates})$ coherence transfers $\rightarrow O(p)$; more memory \\
One lock protects a big table & One lock per bucket / per partition & Contention divided by number of locks; more lock storage \\
Spinning PEs flood the bus & Test-and-test-and-set, back-off & Spins hit in local cache; slight extra latency \\
Barrier after every small step & Merge steps into larger phases; point-to-point flags & Fewer waits; more complex code \\
Ordering between producer and consumer & Bounded queue with flags + fences & Consumer starts as soon as data is ready \\
\bottomrule
\end{tabular}
\end{center}

\begin{trap}[Multi-PE traps]
\begin{itemize}
\item Adding PEs to a memory-bound or lock-bound program makes it no faster, and often slower.
\item Atomics are not free: each one is a coherence transaction on a contended line (Lab 5: CAS retries rose with threads).
\item Padding fixes \emph{false} sharing only; true sharing needs privatisation.
\item Dynamic scheduling is not always better: with equal tasks it adds overhead and loses locality.
\item Speedup must be measured against the best \emph{serial} program, not the parallel program on one PE.
\end{itemize}
\end{trap}

\section{Putting It Together: Scenarios, Dependencies, Checklist}\label{ch:scen}

\subsection{Knock-on effects: what else must change (adjust the whole architecture)}\label{sec:knockon}
Exam answers lose marks for choosing a feature whose prerequisite is missing, or for fixing one layer and creating a bottleneck in the next. For every change you propose, read its row and add the ``also need'' items to your answer.

\small
\begin{longtable}{L{4.6cm}L{5.4cm}L{5.4cm}}
\toprule
\tblhead{\textbf{If you choose\ldots} & \textbf{You also need\ldots} & \textbf{Otherwise\ldots}}\\
\midrule
\endhead
Non-blocking cache (hit-under-miss) & Out-of-order CPU (scoreboard/Tomasulo), multiple memory banks & Blocking CPU stalls anyway; cannot support multiple outstanding misses \\
Hardware speculation & Dynamic branch prediction, reorder buffer (commit in order), recovery/flush & No way to undo wrong-path work; imprecise exceptions \\
Wide issue (superscalar) & Accurate predictor + BTB, more register-file/memory ports, renaming, enough independent instructions & Issue slots idle; clock rate and power suffer \\
SMT & Out-of-order superscalar core, renaming, larger physical register file, per-thread PC/rename/ROB, cache associativity/size increase & Needless on narrow in-order core; shared-resource conflicts \\
Predict-taken static scheme & Branch target known early (BTB/early BTA) & No benefit \\
VLIW & Strong compiler (alias analysis, unrolling, scheduling), no need for binary compatibility & Poor code size/utilisation, lock-step stalls \\
Hardware prefetch & Spare memory bandwidth & Prefetches compete with demand misses \\
Compiler prefetch & Non-blocking cache; issue slots & Wasted instructions; stalls \\
Way prediction / pseudo-associative & Set/column structure; tolerance for variable hit latency (better for L2 than L1 for pseudo) & Pipeline problems when hit time varies \\
Virtual-address cache & PID tags, anti-aliasing (page colouring), protection bits & Flushes on context switch; synonyms \\
More cores / PEs & More memory bandwidth (channels/banks), larger shared LLC, scalable coherence (directory beyond a few cores), less sharing in software & Cores starve on the shared bus; coherence and lock traffic grow with $p$ \\
Wider SIMD / more vector lanes & Proportionally more memory bandwidth, aligned/unit-stride layout, masks for branches & Lanes idle waiting for data or masked off \\
GPU offload & Data resident on GPU or overlapped transfers, coalesced layout (SoA), enough threads for occupancy, little divergence & PCIe transfer dominates; SMs idle \\
Larger OoO window (ROB) & Larger PRF, LQ/SQ, issue queue, more branch checkpoints, accurate predictor, non-blocking cache & Stall just moves to the next full structure (Lab 7) \\
More functional units & Wider issue/wakeup and bypass network, independent instructions of \emph{that} class & No gain unless that class was saturated (Lab 8) \\
Dynamic task scheduling & Cheap concurrent queue, chunk size tuned, locality awareness & Queue contention; lost cache locality \\
Privatised per-PE data & Final reduction step, memory per PE, padding to line size & Wrong result without reduction; false sharing if packed \\
Multithreading for latency hiding & Enough cache capacity/associativity for all threads, more registers & Threads thrash each other's cache \\
\bottomrule
\end{longtable}
\normalsize

\subsection{Connect every concept: the chains examiners look for}\label{sec:connect}
The professor asks you to ``interconnect'' the concepts. Each chain below is one paragraph in an answer: start at the change, follow the arrows, and say what you do at each arrow.
\begin{itemize}
\item \textbf{More PEs} $\rightarrow$ more simultaneous misses $\rightarrow$ shared bandwidth and LLC contention $\rightarrow$ coherence traffic on shared data $\rightarrow$ synchronisation cost $\rightarrow$ Amdahl limit. (Banks/channels, larger LLC, directory, privatise, reduce serial part.)
\item \textbf{Wider issue} $\rightarrow$ needs independent instructions (ILP) $\rightarrow$ limited by dependences and branches $\rightarrow$ better predictor + larger window (ROB/PRF/LQ/SQ) $\rightarrow$ more outstanding misses $\rightarrow$ non-blocking cache/MSHRs $\rightarrow$ power and clock cost.
\item \textbf{DLP hardware} $\rightarrow$ many lanes need many operands per cycle $\rightarrow$ memory bandwidth, unit-stride/coalesced layout $\rightarrow$ branches become masks (divergence) $\rightarrow$ latency hidden by warps (occupancy, registers).
\item \textbf{Deeper pipeline} $\rightarrow$ higher clock $\rightarrow$ larger branch and load penalties $\rightarrow$ dynamic prediction, forwarding $\rightarrow$ more power ($P\propto CV^2f$).
\item \textbf{Multithreading} $\rightarrow$ hides memory latency $\rightarrow$ threads share caches/TLB $\rightarrow$ more capacity/associativity needed $\rightarrow$ single-thread may slow (priority).
\item \textbf{Bigger/more associative caches} $\rightarrow$ fewer misses $\rightarrow$ longer hit time (clock) and more power $\rightarrow$ keep L1 small, grow L2/L3 $\rightarrow$ judge by AMAT, not hit rate (Lab 3).
\item \textbf{Static decisions} $\rightarrow$ need predictable workload $\rightarrow$ cheap hardware, low power; \textbf{dynamic decisions} $\rightarrow$ adapt to data $\rightarrow$ hardware/runtime overhead.
\end{itemize}

\subsection{Lab evidence you can cite}\label{sec:labs}
\small
\begin{longtable}{L{2.4cm}L{7.6cm}L{5.6cm}}
\toprule
\tblhead{\textbf{Lab} & \textbf{What we observed} & \textbf{Use it to argue}}\\
\midrule
\endhead
Lab 1: direct-mapped cache simulator & Sequential access governed by block size; repeated access by whether the working set fits; aliasing addresses gave 100\% miss rate with a tiny working set & Conflict misses need associativity/padding, not a bigger cache \\
Lab 3: three-level hierarchy & Raising L1 associativity fixed a conflict-miss storm, but L2/L3 miss \emph{rates} rose while AMAT and DRAM traffic improved; effects were phase-dependent & Misses move between levels; judge by AMAT/time, not hit rate \\
Lab 4: victim cache, prefetchers, MSHRs & Stride prefetcher cut L1 miss rate 57.4\%$\rightarrow$23.2\% but was slower (0\% timeliness, pollution); victim cache helped only on a conflict pattern (61.6\% fewer cycles, cutoff at $N-$ways entries); MSHRs: no gain for dependent pointer chase, 8$\times$ for 8 independent streams, plateau at 8 & Fewer misses $\ne$ faster; MLP needs \emph{independent} misses; match the optimisation to the access pattern \\
Lab 5: cache coherence & True sharing on a 256\,B counter table: CAS retries rose with threads, IPC collapsed, instructions per event flat; bandwidth, capacity, locks ruled out; per-thread local tables fixed it & Shared writable data destroys scaling; privatise + reduce; rule out alternatives explicitly \\
Lab 6: x86 decoder & 6-instruction/cycle decode ceiling limited compact code, not 32\,B/cycle; microcode-ROM instructions dropped decode to 1 lane; fusion depended on alignment & The front end can be the bottleneck of a wide core \\
Lab 7: OoO allocation & Stalls came from ROB/PRF/IQ/LQ/SQ/checkpoints; enlarging one moved the stall to the next (PRF $\rightarrow$ retire width); window limited by branch accuracy and true dependences & Size structures in balance; bigger windows need independent work \\
Lab 8: OoO scheduler & Ready $\ne$ issued; time set by the longest dependence chain unless a unit class saturates; 48/96-entry scheduler, 8-wide issue, 4 ALUs: same 123 cycles; fixing loads moved the limit to the multiplier & Add resources only of the saturated class; the bottleneck moves \\
Lab 10: producer--consumer bypass & One recurrence: 0.27 IPC vs.\ $\approx$3 for independent work (11$\times$); about 11 interleaved chains needed to fill the ports & Break dependence chains (multiple accumulators, threads); FUs cannot fix a recurrence \\
\bottomrule
\end{longtable}
\normalsize

\subsection{Worked scenarios (model answers)}

\subsubsection*{Scenarios 1--5 --- Reference designs by machine class (use when the question names the class)}
{\footnotesize
\begin{longtable}{L{1.45cm}L{2.8cm}L{2.8cm}L{2.8cm}L{2.8cm}L{2.8cm}}
\toprule
\tblhead{ & \textbf{1. Phone / PMD} & \textbf{2. Web/DB server} & \textbf{3. Desktop, 1 thread} & \textbf{4. Scientific FP} & \textbf{5. Embedded}}\\
\midrule
\endhead
Goal & Energy per task, real-time & Throughput, availability, cost/request & Single-thread latency, price-performance & FP throughput, huge arrays & Price \\
Core & 2-way in-order or modest OoO, 5--8 stages & Moderate 4-issue OoO, or simpler fine-grained MT cores & 4-wide Tomasulo + ROB speculative & Pipelined FP add/mult; pipelined or duplicated divider; unrolling + software pipelining (VLIW OK) or Tomasulo if latencies vary & 5-stage in-order, forwarding, branch resolved in ID \\
Branch & Small 2-bit BHT + small BTB + RAS & Good dynamic predictor & Two-level/gshare + BTB + RAS & BTFNT or 2-bit (loops) & Static (not-taken/BTFNT) + 1 delay slot \\
L1 & 32--64\,KB, 2-way + way prediction & Private & 32--64\,KB, low associativity, 64\,B & Blocking, interchange, fusion; stream prefetch; non-blocking & Small direct-mapped I/D \\
L2 / L3 & 256\,KB L2, no L3 & Large shared L3 (16--64\,MB) & 256\,KB L2 + 8--32\,MB L3 & 64\,B blocks, prefetch & None or small \\
Memory & Interleaved low-power DRAM & 32--256\,GB, banked, non-blocking, ECC & Non-blocking + HW prefetch & Wide / heavily interleaved & Simple \\
Threads, cores & Several small cores, no SMT & Many cores + 2--4 SMT/FGMT threads & 2--4 cores, optional SMT with priority & Multicore + SIMD/vector/GPU & One core \\
Power / other & DVFS, power gating & Coherence-aware software & --- & Roofline check first & VLIW only for fixed numeric firmware \\
Rejected & Wide speculative core, trace cache, big L3, VLIW & Very wide single core, coarse MT only, deep speculation & VLIW, coarse MT, very wide issue & Big predictor, heavy speculation, SMT (FP gains least) & Tomasulo+ROB, BTB, trace/victim cache \\
\bottomrule
\end{longtable}}
\textbf{Key reasons to quote:} ILP techniques are energy-inefficient (PMD); DB/web are the best multicore beneficiaries and ILP of one thread leaves most units idle (server); 3--6 issue is the practical ILP sweet spot, a deep pipe needs a strong predictor (desktop); unrolling took the FP loop from 9 to 3.5 cycles/iteration, interchange cut misses by 93.75\% (scientific); static techniques need no predictor memory and software cache optimisations cost no hardware (embedded).
\subsubsection*{Exam-style model answers: a given architecture that is not working ideally}
Scenarios 6--8 follow the exact format the professor described (architecture + workload + non-ideal behaviour) and the answer layout of \S\ref{sec:8step}. Numbers are constructed for practice.

\subsubsection*{Scenario 6 --- Multicore analytics that gets slower with more cores}
\textbf{Given.} 8-core shared-memory chip, private 32\,KB L1 per core, shared 8\,MB L2, MESI snooping bus, 64\,B lines. Workload: 80\,M independent events streamed from memory (16\,B each); for each event a core parses it and atomically increments one of 16 shared 8-byte counters. \textbf{Symptom:} 1 core 10\,s; 4 cores 6\,s; 8 cores 7.7\,s (slower than 4).

\textbf{1. Assumptions.} A1: events cost the same. A2: input read once from DRAM. A3: counters are one contiguous 128\,B array (2 cache lines).

\textbf{2. Workload analysis.} Per-event work is independent (TLP/DLP, parallel fraction $\approx1$). The counter update is a \emph{reduction} (associative, commutative), so the order of increments does not matter: it can be done privately and combined later. Memory: 1.28\,GB streamed once (spatial locality only); counters 128\,B (tiny, fit anywhere).

\textbf{3. Why not ideal.} Every increment needs the counter line in M state. With 8 writers on 2 lines, almost every increment is a coherence miss: invalidate other copies, transfer the line over the bus, retry the atomic. Updates to one line are serialised, so the effective serial fraction \emph{grows with} the core count (true sharing, plus false sharing among the 8 counters in one line). Ruled out: capacity (128\,B), bandwidth (1.28\,GB in 7.7\,s $\approx$0.17\,GB/s, far below DRAM bandwidth), locks (none used). Same signature as Lab 5 (CAS retries up, IPC down, instructions per event flat).

\textbf{4. Proposed architecture.}
{\small
\begin{longtable}{L{2.6cm}L{6.0cm}L{1.3cm}L{5.6cm}}
\toprule
\tblhead{\textbf{Layer} & \textbf{Change} & \textbf{S/D} & \textbf{Why}}\\
\midrule
\endhead
Distribution & Block partition: core $k$ takes events $[k\cdot10\text{M},(k{+}1)\cdot10\text{M})$ & Static & Equal costs, contiguous data (prefetch-friendly), zero scheduling overhead \\
Aggregation & Private 16-counter array per core, aligned and padded to its own lines & Static & Updates hit in own L1 in M/E state: no coherence traffic in the hot loop \\
Synchronisation & One barrier, then tree reduction (3 steps for 8 cores) & Static & 80\,M atomics $\rightarrow$ 1 barrier + 7 merges \\
Caches/coherence & Unchanged MESI; E state makes first write silent & --- & Sharing removed, so the protocol cost disappears \\
Memory & Hardware stream prefetch for the input & Dynamic & Regular sequential pattern; keeps the bus busy \\
Core & Unchanged (no wider core, no SMT) & --- & Not ILP-bound; problem was sharing \\
\bottomrule
\end{longtable}}
\textbf{5. Why better.} Coherence transfers drop from $O(\text{events})$ to $O(p)$. Rejected: a lock (serialises everything, worse); padding alone (fixes false sharing, but all cores still write the \emph{same} counters: true sharing remains); faster interconnect or directory (each transfer cheaper, but still 80\,M of them); more cores (more contention).
\textbf{6. Demerits / next bottleneck.} $8\times128$\,B extra memory, a merge step, code change. Expect near 8$\times$ ($\approx$1.3\,s) until input bandwidth binds ($1.28\,\text{GB}/1.3\,\text{s}\approx1$\,GB/s, still fine).

\subsubsection*{Scenario 7 --- GPU image filter at 8\% of peak}
\textbf{Given.} GPU with 40 SMs, 32-thread warps, PCIe 16\,GB/s to the host. Workload: per-frame filter on a 3840$\times$2160 RGBA image (4\,B/pixel, $\approx$33\,MB, row-major). Kernel: thread $i$ processes row $i$, walking across its columns. 30\% of pixels (randomly located ``edge'' pixels) need a 50-instruction path; others need 5. Each frame is copied in, processed, copied out. \textbf{Symptom:} 8\% of peak; kernel 6\,ms, copies $\approx$2\,ms each way.

\textbf{1. Assumptions.} A1: up to 48 resident warps per SM. A2: the filter reads a $3\times3$ neighbourhood.

\textbf{2. Workload analysis.} Pure DLP: each output pixel is independent (no order between pixels); neighbourhoods overlap (reuse); the branch is data-dependent.

\textbf{3. Why not ideal (four causes).} (a) \emph{Occupancy}: 2160 threads $=68$ warps for 40 SMs ($<2$ per SM, vs.\ 48 possible), so memory latency cannot be hidden by switching warps. (b) \emph{Uncoalesced access}: at each step the 32 threads of a warp read 32 different rows, 15\,KB apart: up to 32 transactions instead of one. (c) \emph{Divergence}: $P(\text{no edge pixel in a warp})=0.7^{32}\approx0$, so every warp runs both paths: 55 instructions vs.\ 18.5 useful (34\% SIMD efficiency, \S\ref{sec:divergence}). (d) \emph{Transfer}: 4\,ms of copies per 6\,ms kernel.

\textbf{4. Proposed architecture.}
{\small
\begin{longtable}{L{2.6cm}L{6.0cm}L{1.3cm}L{5.6cm}}
\toprule
\tblhead{\textbf{Layer} & \textbf{Change} & \textbf{S/D} & \textbf{Why}}\\
\midrule
\endhead
Work mapping & One thread per pixel, \texttt{threadIdx.x} along the row, $16\times16$ blocks & Static & 8.3\,M threads $\Rightarrow$ full occupancy; neighbours in a warp read neighbouring addresses (coalesced) \\
On-chip memory & Each block loads its tile + 1-pixel halo into shared memory once & Static & $3\times3$ neighbourhood reused 9$\times$ from scratchpad: higher arithmetic intensity \\
Divergence & Pass 1 classifies pixels and writes a compacted list of edge pixels (prefix sum); pass 2 runs the 50-instruction path on full warps & Dynamic & Decision depends on data, known only at run time; warps become uniform \\
Latency hiding & Hardware warp scheduler with many warps & Dynamic & Fine-grained MT covers DRAM latency \\
Host link & Double buffering: copy frame $k{+}1$ while computing frame $k$; keep intermediate data on GPU & Static & Copies overlap compute instead of adding to it \\
Sync & Barrier only inside a block (after tile load); atomics only for the list counter, aggregated per warp & Static & No cross-block barrier needed \\
\bottomrule
\end{longtable}}
\textbf{5. Why better.} Rejected: a bigger GPU (more SMs, same waste per warp); predication alone (still pays both paths); an 8-core CPU with 8-wide SIMD (64 lanes vs.\ thousands; but it avoids PCIe, so it is the right choice only for small images).
\textbf{6. Demerits / next bottleneck.} Extra classification pass and list memory; one frame of added latency from double buffering; the kernel becomes memory-bandwidth bound (check with the roofline, \S\ref{sec:roofline}).

\subsubsection*{Scenario 8 --- Wide core, low IPC: the bottleneck moves}
\textbf{Given.} 4-wide OoO core at 3\,GHz, 2 pipelined FMA units (latency 4), 32\,KB L1, 1\,MB L2, DRAM 20\,GB/s. Workload: \texttt{for i: sum += a[i]*b[i]} over two 64\,M-element double arrays (512\,MB each). \textbf{Symptom:} IPC $\approx$1.25 instead of $\approx$4. The vendor proposes an 8-wide core.

\textbf{Analysis and diagnosis.} 5 instructions per iteration (2 loads, FMA, increment, branch). Each FMA needs the previous \texttt{sum}: a loop-carried chain of 4 cycles per iteration, so $5/4=1.25$ IPC whatever the width (Lab 8, Lab 10). The 8-wide core changes nothing.

\textbf{Fix and where the bottleneck moves.}
(1) \emph{Static, compiler}: unroll and use $4\times2=8$ independent partial sums (latency $\times$ units), combined at the end (FP reassociation slightly changes rounding: state it). Now 2 FMA/cycle $\Rightarrow$ 6\,G iterations/s possible.
(2) \emph{Memory check}: each iteration moves 16\,B; 20\,GB/s allows only 1.25\,G iterations/s. Before: 0.75\,G/s. So the gain is $1.25/0.75\approx1.7\times$, not 8$\times$: the workload is now \textbf{bandwidth-bound} (AI $=2/16=0.125$ FLOP/B). SIMD would add nothing.
(3) \emph{Whole architecture}: more memory channels/banks, hardware stream prefetch (dynamic, regular pattern) to reach peak bandwidth; split across cores/sockets with their own memory, each with a private partial sum and a final reduction; caches do not help (no reuse: 1\,GB read once), unless the program can be fused with other passes over the same data to raise AI.
\textbf{Why better:} targets the actual binding terms in order (recurrence, then bandwidth). Rejected: 8-wide core (neither term), bigger cache (no reuse), SMT alone (second thread breaks the chain but shares the same bandwidth).

\subsubsection*{Scenario 9 --- Numeric cache comparison (practice, constructed numbers)}
\textbf{Problem.} L1 miss penalty 40\,cc, hit time 1\,cc. Option A: 16\,KB direct-mapped, miss rate 3.0\%. Option B: 16\,KB 2-way, miss rate 2.6\%, clock stretched by 10\% (slide: 2-way $\times1.10$). Option C: option A plus a 4-entry victim cache removing 40\% of misses (slide range 20--95\%). Which is best?

\begin{ex}[Solution (AMAT in direct-mapped cycles)]
A: $1+0.030\times40=2.20$. B: $1.10+0.026\times40=2.14$ --- only 2.7\% better \emph{in AMAT}, but the 10\% longer clock also slows \emph{every non-memory instruction}, so B is a net loss. C: $1+0.030\times0.6\times40=1.72$ (assuming the victim probe hides inside the normal miss path, per the slide). \textbf{Choose C}: keeps the direct-mapped hit time (clock) and removes conflict misses cheaply. Quote the slide table: at 16\,KB, 1-way has the lowest AMAT.
\end{ex}

\subsubsection*{Scenario 10 --- Numeric branch-handling comparison (practice, constructed numbers)}
\textbf{Problem.} Deep pipeline (R4000-like): branch penalty 3 cycles, 20\% of instructions are branches, 60\% of branches are taken, base CPI 1.0. Compare (a) stall always, (b) predict-not-taken, (c) 2-bit BHT+BTB 90\% accurate, (d) correlating/two-level predictor 95\% accurate.

\begin{ex}[Solution]
\begin{center}
\small
\begin{tabular}{L{6cm}C{4.2cm}C{2.4cm}}
\toprule
\tblhead{\textbf{Scheme} & \textbf{Extra CPI} & \textbf{CPI}}\\
\midrule
(a) Stall always & $0.20\times3=0.60$ & 1.60 \\
(b) Predict not taken (wrong when taken) & $0.20\times0.60\times3=0.36$ & 1.36 \\
(c) 2-bit BHT + BTB, 90\% & $0.20\times0.10\times3=0.06$ & 1.06 \\
(d) Two-level, 95\% & $0.20\times0.05\times3=0.03$ & 1.03 \\
\bottomrule
\end{tabular}
\end{center}
Speedup of (c) over (a): $1.60/1.06\approx1.5$; (d) over (c): $\approx1.03$. \textbf{Conclusion:} in a deep pipe the dynamic predictor is mandatory; the extra complexity of (d) is justified only if area allows (diminishing returns), or if branch frequency is higher or the pipeline deeper.
\end{ex}

\subsection{Conflict resolution: when two constraints disagree}
\begin{center}
\small
\begin{tabular}{L{5.2cm}L{10.2cm}}
\toprule
\tblhead{\textbf{Conflict} & \textbf{Resolution}}\\
\midrule
High clock vs.\ low CPI & Deeper pipeline raises clock but adds load/branch delay; compensate with prediction and forwarding. Decide by the \emph{product} IC$\times$CPI$\times T_{clk}$. \\
Low miss rate vs.\ fast hit & Large/associative cache cuts misses but stretches the clock; use small DM L1 + big associative L2; victim cache/way prediction as compromises. \\
Throughput vs.\ single-thread latency & Fine-grained MT/SMT raise throughput but slow a thread; add priority or fewer threads; or heterogeneous cores. \\
Performance vs.\ power & ILP/speculation/prefetch cost energy; prefer TLP/multicore and DVFS; energy per task is the right metric. \\
Compiler simplicity vs.\ hardware simplicity & Dynamic scheduling = simple compiler, complex hardware; VLIW = opposite. Choose by compatibility and cost. \\
Bandwidth vs.\ latency & Buy bandwidth (wide/interleaved memory, banks); hide latency (OoO, non-blocking, MT, prefetch). \\
Compatibility vs.\ simplicity & Compatibility $\Rightarrow$ superscalar; fixed product $\Rightarrow$ VLIW. \\
\bottomrule
\end{tabular}
\end{center}

\subsection{Final ten-minute exam checklist}\label{sec:checklist}
\begin{enumerate}[leftmargin=*]
\item Did I write numbered \textbf{assumptions} for everything the question did not give?
\item Did I explain \textbf{why the problem occurs} (mechanism + a number + one alternative ruled out), before proposing anything?
\item Did I state whether \textbf{parallelism is present}, which kind, the \textbf{order of execution} (dependences, serial part, reductions), and whether \textbf{the memory hierarchy will work}?
\item Did I say how the work is \textbf{distributed over the PEs}, with timing, and what \textbf{communication and synchronisation} it needs?
\item Did I mark each choice \textbf{static or dynamic} and say why?
\item Did I touch \textbf{every layer} the question involves and check knock-on effects (\S\ref{sec:knockon})?
\item Did I give \textbf{benefits and demerits} and say \textbf{why my solution beats at least two alternatives}? Did I cite a \textbf{lab} finding?
\item Did I state the machine class and its priority (energy / cost / throughput / availability / single-thread)?
\item For \emph{every} block (core, predictor, L1, L2/L3, memory, threads) did I give a choice \emph{and} a reason tied to a constraint?
\item Did I give at least one number (AMAT, CPI, stalls/instr, speedup)?
\item Did I check feature dependencies (Section above)? (OoO $\Rightarrow$ non-blocking; speculation $\Rightarrow$ ROB+predictor; SMT $\Rightarrow$ OoO.)
\item Did I state costs (area, power, clock, complexity) of each choice?
\item Did I list rejected alternatives and why?
\item If the scenario hints at compatibility, did I avoid VLIW (or justify)?
\end{enumerate}

\appendix
\section{Formula and Number Sheet}\label{app:formulas}
Numbers worth quoting that are not repeated in the sections above. $\dagger$ = standard formula not spelled out in the slide text.

\subsection{Performance, power, cost, dependability}
\begin{center}
\small
\begin{tabular}{L{5.4cm}L{10cm}}
\toprule
\tblhead{\textbf{Item} & \textbf{Formula}}\\
\midrule
CPU time & $\text{IC}\times\CPI\times T_{clk}$ \\
Average CPI & $\sum_i(\text{IC}_i\times\CPI_i)/\text{IC}$ \\
Speedup & $\text{Time}_{old}/\text{Time}_{new}$ \\
Amdahl$^\dagger$ & $1/[(1-f)+f/s]$ \\
Dynamic energy & $\tfrac12 C V^2$ \\
Dynamic power & $\tfrac12 C V^2 f$ \\
Static power & $I_{static}\times V$ (reduce by power gating) \\
Die yield$^\dagger$ & $\text{wafer yield}\times(1+\text{defects/area}\times\text{die area})^{-N}$, $N=11.5$--$15.5$ \\
MTBF; availability & $\text{MTBF}=\text{MTTF}+\text{MTTR}$; $\text{Availability}=\text{MTTF}/\text{MTBF}$ \\
Price vs.\ volume & $-10\%$ per doubling of volume \\
\bottomrule
\end{tabular}
\end{center}

\subsection{Memory}
\begin{center}
\small
\begin{tabular}{L{6.2cm}L{9.2cm}}
\toprule
\tblhead{\textbf{Item} & \textbf{Formula / value}}\\
\midrule
AMAT & $t_{hit}+m\times P$ \\
Two-level AMAT & $t_{L1}+m_{L1}(t_{L2}+m^{loc}_{L2}P_{L2})$; global $m_{L2}=m_{L1}m^{loc}_{L2}$ \\
Stalls/instruction (two-level) & $\text{misses}_{L1}/\text{instr}\times t_{L2}+\text{misses}_{L2}/\text{instr}\times P_{L2}$ \\
Slide example & $1+0.04(10+0.5\times100)=3.4$; stalls $(3.4-1)\times1.5=3.6$ \\
L2 associativity example & DM $10+0.25\times100=35$; 2-way $10+0.20\times100=30$ \\
CPI with memory & $\CPI_{base}+\text{refs/instr}\times m\times P$ (e.g.\ $2+1.2\times0.03\times64=4.3$) \\
Block-size penalty model & 80 cc + 2 cc per 16 B: 82/84/88/96/112 for 16/32/64/128/256 B \\
Best block per cache (AMAT) & 4K: 32 B; 16K, 64K, 256K: 64 B \\
Associativity clock factors & $\times1.10$ (2-way), $\times1.12$ (4-way), $\times1.14$ (8-way) \\
2:1 rule & 2-way of size $N/2\approx$ DM of size $N$ (caches $<128$\,KB) \\
Victim cache & 20--95\% miss reduction, 4\,KB DM, 1--5 entries \\
Way prediction accuracy & $>85\%$ (Alpha 21264, 2-way I-cache) \\
Prefetch AMAT & $t+m\,p\,t_p+m(1-p)P$; example $2.41525$ \\
Stream buffers & ISB 1/4/16 blocks catch 15--25\%/50\%/72\% of I-misses \\
Loop interchange & $1{,}048{,}576\to65{,}536$ (per array); $-93.75\%$ \\
Loop fusion & $262{,}144\to196{,}608$ ($-25\%$) \\
Blocking & $2N^3+N^2\to2N^3/B+N^2$; 3 tiles of $B^2$ words must fit \\
Interleaving & bank = addr mod $\#$banks; banks $\ge$ cycles to access a word \\
Memory-width CPI example & 1-word 4.30; 2-word: simple 5.07 / interleaved 3.63 / wide 3.54; 4-word: 5.69 / 3.09 / 3.84 \\
TLB & 128--256 entries typical; fully associative on high-end \\
Page size / VA split & 4--8\,KB; 32-bit VA = 20-bit VPN + 12-bit offset \\
DRAM refresh & $\sim$8\,ms, $\sim$1\% of time \\
\bottomrule
\end{tabular}
\end{center}

\subsection{Threads and multicore}
\begin{center}
\small
\begin{tabular}{L{6.2cm}L{9.2cm}}
\toprule
\tblhead{\textbf{Item} & \textbf{Value}}\\
\midrule
Pentium 4 SMT & 1.01 (int rate), 1.07 (FP rate); pairs 0.90--1.58, avg 1.20 \\
Power5 SMT & 1.23 (int rate), 1.16 (FP rate); 0.89--1.41 for two copies \\
Power5 core size & $+24\%$ vs.\ Power4; physical regs 152$\to$240; L2 1.44$\to$1.92\,MB \\
SMT threads per core & 2, 4, or sometimes 8 (Power5: 2) \\
Hardware thread switch & vs.\ 100s--1000s of clocks for a process switch \\
General SMT gain & up to 30\% \\
\bottomrule
\end{tabular}
\end{center}

\subsection{Parallel execution, DLP and communication}
\begin{center}
\small
\begin{tabular}{L{5.0cm}L{10.4cm}}
\toprule
\tblhead{\textbf{Item} & \textbf{Formula}}\\
\midrule
Amdahl / Gustafson & $S=1/[(1-f)+f/p]$ (limit $1/(1-f)$); $S=(1-f)+fp$ \\
Efficiency; imbalance & $E=S/p$; imbalance $=\max_i W_i/\text{avg}_i W_i$ \\
Parallel time & $T_{par}=T_{serial}+\max_i(W_i/\text{speed}_i)+T_{comm}+T_{sync}$ \\
Work--span & $T_p\ge\max(T_1/p,\,T_\infty)$; parallelism $=T_1/T_\infty$ \\
Tree reduction & $\log_2 p$ combine steps \\
Message time & latency $+$ size/bandwidth \\
Grid boundary per PE ($N\times N$) & strips $\approx2N$; blocks $\approx4N/\sqrt{p}$ \\
Roofline & $\min(\text{peak FLOP/s},\ \text{AI}\times\text{BW})$, AI $=$ FLOPs/DRAM bytes \\
Vector time & convoys $\times\ n/L$ + start-up; strip-mine when $n>$ MVL \\
Bank conflicts & unit stride needs banks $\ge$ busy time; stride $s$ conflicts if banks$/\gcd(s,\text{banks})<$ busy time \\
SIMT efficiency & active lanes / 32; divergent warp time $\approx T_{if}+T_{else}$ \\
Occupancy & resident warps / max warps per SM (limited by registers, shared memory, block size) \\
Independent accumulators to hide latency & latency $\times$ number of pipelined units (e.g.\ $4\times2=8$) \\
\bottomrule
\end{tabular}
\end{center}

\subsection{Typical values for writing assumptions}\label{sec:typical}
Use these only when the question does not give a value, and label them as assumptions.\bs
\begin{center}
\small
\begin{tabular}{L{5.0cm}L{10.4cm}}
\toprule
\tblhead{\textbf{Item} & \textbf{Typical value}}\\
\midrule
Cache line & 64\,B \\
L1 / L2 / L3 & 32--64\,KB, 4--8-way, 1--4 cycles / 256\,KB--2\,MB, 8--16-way, $\approx$10--15 cycles / 8--64\,MB shared, $\approx$30--50 cycles \\
DRAM latency; bandwidth & 60--100\,ns ($\approx$200--300 cycles at 3\,GHz); 20--50\,GB/s per socket (CPU), several hundred GB/s to TB/s (GPU) \\
Branch misprediction penalty & 10--20 cycles on a deep OoO core; 1--3 on a short in-order pipe \\
Core-to-core line transfer & tens of cycles on-chip; $\sim$100+ cycles across sockets \\
GPU & warp 32 threads; tens of resident warps per SM; PCIe 16--32\,GB/s \\
Issue width & 2 (embedded), 4--6 (desktop/server OoO) \\
\bottomrule
\end{tabular}
\end{center}

\section{Known Inconsistencies in the Slides}\label{app:src}
\textbf{R4000 load delay:} one slide title says ``x1'', the table shows \emph{two} stall cycles and the older deck says ``TWO cycle load latency'': use \textbf{2}.
\textbf{Memory-width CPI example:} buses are labelled ``64-bit''/``128-bit'' but the arithmetic treats them as 1-word/2-word; quote the CPI values (4.30, 5.07, 3.63, 3.54, 5.69, 3.09, 3.84), not the labels.
\textbf{Blocking numbers} ($\approx1.07$ billion $\to\approx2.1$ million misses) come from a simplified model; present them as illustrative.


\end{document}
